\documentclass[11pt,oneside,openany]{book}
\usepackage[T1]{fontenc}
\usepackage{lmodern,microtype}
\usepackage[a4paper,margin=28mm]{geometry}
\usepackage{amsmath,amssymb,amsthm,mathtools}
\usepackage[hidelinks,bookmarksopen,bookmarksnumbered]{hyperref}
\hypersetup{pdftitle={Extremal Methods, Networks, and Matroid Optimization},
  pdfauthor={MathGenome},pdfsubject={Unified mathematical notes, P2-05}}
\allowdisplaybreaks[1]
\setlength{\emergencystretch}{2em}

\theoremstyle{plain}
\newtheorem{theorem}{Theorem}[chapter]
\newtheorem{lemma}[theorem]{Lemma}
\newtheorem{proposition}[theorem]{Proposition}
\newtheorem{corollary}[theorem]{Corollary}
\theoremstyle{definition}
\newtheorem{definition}[theorem]{Definition}
\newtheorem{example}[theorem]{Example}
\theoremstyle{remark}
\newtheorem{remark}[theorem]{Remark}

\begin{document}
\frontmatter
\hypersetup{pageanchor=false}
\begin{titlepage}
\thispagestyle{empty}
\vspace*{27mm}
{\large MathGenome \enspace / \enspace P2-05\par}
\vspace{8mm}
\hrule height 0.8pt
\vspace{11mm}
{\Huge\bfseries Extremal Methods, Networks,\\[3pt]
and Matroid Optimization\par}
\vspace{12mm}
{\large Unified mathematical notes\par}
\vspace{6mm}
{\normalsize Extremal graphs, set families, and Ramsey theory\\
The probabilistic method, concentration, and the local lemma\\
Linear programming, network flows, assignment, and matroids\par}
\vfill
{\large 8 November 2026\par}
\end{titlepage}
\hypersetup{pageanchor=true}
\begingroup
\small
\setlength{\baselineskip}{10.6pt}
\tableofcontents
\endgroup

\chapter*{Scope and conventions}
\addcontentsline{toc}{chapter}{Scope and conventions}
The first three chapters study finite extremal bounds, shadows of
set families, and Ramsey theory. Chapters 4--6 develop the probabilistic
method from existence and alteration to moment bounds, concentration,
and the local lemma. The last three chapters develop dual certificates,
integrality, network and assignment algorithms, and matroid
optimization. The companion result map lists fifty writing groups,
including twenty-two original contracts, with statement and proof
locations.

Extremal graphs are finite simple undirected graphs unless an infinite
ground set is explicitly specified. Directed networks have labelled
arcs and use net outflow for divergence. Matroids have a finite
specified ground set. A finite probability model includes its
outcomes and probabilities, and statements on general probability
spaces specify the required measurability and integrability. Empty
sums and products are \(0\) and \(1\), respectively.

The optimization chapters use the proved finite-dimensional
separation, Farkas, and linear programming results in \cite{p106},
and the matching, graph, network, and finite matroid results in
\cite{p204}. Exact references are given where these results are used.
The probability arguments use elementary measure and integration
from \cite{d11}. Textbook citations identify the source or a parallel
treatment. The proofs below give the mathematical arguments for the
new results, including the finite transport and conditioning steps
that are needed to apply them.

\mainmatter
\chapter{Extremal Graphs}

All graphs in this chapter are finite, simple, and undirected. A copy
means a subgraph copy, so it may have extra edges in the ambient graph.
Write \(e(G)\) for the number of edges of \(G\).

\section{Clique extremal numbers}

\begin{theorem}[Mantel and Tur\'an]\label{p205:W01}
For integers \(n\ge 0\) and \(r\ge 1\), let \(T_r(n)\) be the complete
\(r\)-partite graph on \(n\) vertices with part sizes differing by at
most one. Parts may be empty. Every \(K_{r+1}\)-free graph on \(n\)
vertices has at most \(e(T_r(n))\) edges. Equality holds exactly for
graphs isomorphic to \(T_r(n)\). For \(n=qr+b\) with \(0\le b<r\),
\[
 e(T_r(n))=\binom n2-(r-b)\binom q2-b\binom{q+1}2.
\]
In particular, a triangle-free \(n\)-vertex graph has at most
\(\lfloor n^2/4\rfloor\) edges.
\end{theorem}
\begin{proof}
The case \(n=0\) is immediate. Choose a \(K_{r+1}\)-free graph with the
maximum number of edges. If \(u,v\) are distinct nonadjacent vertices,
replace the neighborhood of \(u\) by that of \(v\). This creates no
\(K_{r+1}\). Such a clique through \(u\) cannot use \(v\), and substituting
\(v\) for \(u\) would give a clique in the old graph. The edge change is
\(\deg(v)-\deg(u)\). Reversing the roles of \(u,v\) shows that
\(\deg(u)=\deg(v)\).

Suppose \(u,w\) are adjacent and each is nonadjacent to a third vertex
\(v\). Their degrees both equal \(\deg(v)\). Delete their incident
edges, and connect each of \(u,w\) to the old neighborhood of \(v\).
The two replaced vertices are nonadjacent and neither is adjacent to
\(v\). Any clique through either can therefore again be transferred
to an old clique through \(v\). The edge gain is
\(2\deg(v)-(\deg(u)+\deg(w)-1)=1\), a contradiction. Thus the relation
of equality or nonadjacency is transitive and is an equivalence
relation. Its classes are the parts of a complete multipartite graph.
There are at most \(r\) nonempty parts because one point from each part
forms a clique.

Pad the list of part sizes to \(r\) with zeros. For part sizes \(a\ge b+2\),
moving one point from the larger to the smaller reduces the missing
edge count by
\(\binom a2+\binom b2-\binom{a-1}2-\binom{b+1}2=a-b-1>0\).
The resulting graph is still \(K_{r+1}\)-free, so all sizes in a
maximizer differ by at most one. Such a graph exists and is
\(K_{r+1}\)-free, and its part sizes are uniquely specified up to
permutation. Subtracting the within-part pairs proves the formula.
For \(r=2\), the two sizes are \(\lfloor n/2\rfloor,\lceil n/2\rceil\),
with product \(\lfloor n^2/4\rfloor\).
\end{proof}

\section{Forbidden complete bipartite graphs}

\begin{theorem}[K\H{o}v\'ari--S\'os--Tur\'an]\label{p205:W02}
Let \(1\le s\le t\) be integers. If \(G\) has \(n\) vertices and no
copy of \(K_{s,t}\), then
\[
  2e(G)\le (s-1)n+(t-1)^{1/s}n^{\,2-1/s}.
\]
For fixed \(s,t\) this gives the upper bound \(O(n^{\,2-1/s})\).
\end{theorem}
\begin{proof}
The empty graph gives the \(n=0\) case. The case \(s=1\) follows
from the maximum degree bound \(t-1\), which also handles \(t=1\).
Assume \(n\ge1\) and \(s\ge2\). Count pairs \((v,S)\) where \(S\)
consists of \(s\) neighbors of \(v\). Every \(s\)-set \(S\) has at
most \(t-1\) common neighbors. A common neighbor cannot belong to
\(S\), since there are no loops. Thus
\[
  \sum_v\binom{\deg(v)}s\le(t-1)\binom ns.
\]
For a nonnegative integer \(d\) one has
\(s!\binom ds\ge \max\{d-s+1,0\}^s\). For \(d\ge s\) use the \(s\)
factors of the falling factorial, and for \(d<s\) both sides are
zero. Set \(y_v=\max\{\deg(v)-s+1,0\}\) and
\(\bar y=n^{-1}\sum_v y_v\). For \(x\ge0\) the inequality
\(x^s\ge \bar y^s+s\bar y^{s-1}(x-\bar y)\) follows by integrating
the monotone derivative of \(x^s\) between \(x\) and \(\bar y\).
Summing at \(x=y_v\) gives
\[
 n\bar y^s\le\sum_v y_v^s\le s!(t-1)\binom ns\le (t-1)n^s.
\]
Since \(2e(G)-(s-1)n\le \sum_v y_v\), taking the positive \(s\)-th
root yields the stated bound.
\end{proof}

\section{Dense graphs and multipartite subgraphs}

For \(r,t\ge1\), write \(K_r(t)\) for the complete \(r\)-partite
graph with \(t\) vertices in every part.

\begin{lemma}[Dense families of transversal tuples]\label{p205:transversal-density}
Let \(r,t\ge1\) and \(0<c\le1\) be fixed. For all sufficiently large
\(n\), suppose \(V_1,\dots,V_r\) are finite sets of size at most \(n\)
and \(\mathcal E\subseteq V_1\times\cdots\times V_r\) has at least
\(cn^r\) tuples. There exist \(t\)-subsets \(T_i\subseteq V_i\)
such that \(T_1\times\cdots\times T_r\subseteq\mathcal E\).
\end{lemma}
\begin{proof}
When \(r=1\), choose \(t\) elements among the at least \(cn\) allowed
points. Suppose \(r>1\), and for a prefix \(u\) of length \(r-1\)
let \(N(u)\subseteq V_r\) be the set of its extensions. Prefixes with
\(|N(u)|<cn/2\) contribute at most \(cn^r/2\) tuples in total. Since
each prefix contributes at most \(n\), at least \(cn^{r-1}/2\)
prefixes have \(|N(u)|\ge cn/2\). For large enough \(n\), say \(cn/2\ge2t\),
each such prefix has at least \((cn/4)^t/t!\) distinct \(t\)-subsets
of \(N(u)\), by the falling factorial formula. Count pairs of a
prefix and such a subset. There are at most \(n^t/t!\) subsets on
the last part, so one \(t\)-set \(T_r\) extends at least
\((c/2)(c/4)^t n^{r-1}\) prefixes. Apply the induction hypothesis
to the family of these prefixes. The resulting product can be
extended by every point of \(T_r\).
\end{proof}

\begin{theorem}[Clique blow-ups above the Tur\'an threshold]\label{p205:W03}
Fix integers \(r\ge2,\ t\ge1\) and \(\varepsilon>0\). For all
sufficiently large \(n\), every \(n\)-vertex graph \(G\) with
\[
  e(G)\ge \left(1-\frac1{r-1}+\varepsilon\right)\binom n2
\]
contains \(K_r(t)\).
\end{theorem}
\begin{proof}
If \(1-1/(r-1)+\varepsilon>1\), the hypothesis is impossible for
\(n\ge2\). Otherwise put \(d=1-1/(r-1)\). Choose a fixed integer
\(m\ge r\) large enough that \(d/(m-1)\le\varepsilon/2\).
For every partition into \(r-1\) parts with sizes \(a_i\), the number
of crossing pairs is \((m^2-\sum_i a_i^2)/2\). The inequality
\(\sum_i(a_i-m/(r-1))^2\ge0\) gives the bound \(d m^2/2\).
Consequently Theorem~\ref{p205:W01} gives
\[
  e(T_{r-1}(m))\le(d+\varepsilon/2)\binom m2.
\]

For \(n\ge m\), choose an \(m\)-set uniformly. Each edge is included
with probability \(\binom m2/\binom n2\), so the expected number of
sampled edges is at least \((d+\varepsilon)\binom m2\). Write \(q\)
for the proportion of samples containing an \(r\)-clique. A sample
outside this proportion has at most \((d+\varepsilon/2)\binom m2\)
edges, and every sample has at most \(\binom m2\). Hence the
expected number is at most \((d+\varepsilon/2+q)\binom m2\) and
\(q\ge\varepsilon/2\). Each \(r\)-clique lies in
\(\binom{n-r}{m-r}\) samples, so the number of \(r\)-cliques is at
least
\[
  \frac{\varepsilon}{2}\frac{\binom nm}{\binom{n-r}{m-r}}
    =\frac{\varepsilon}{2}\frac{\binom nr}{\binom mr}.
\]
For \(n\ge2r\), the last expression is at least a fixed positive
multiple of \(n^r\).

Independently label each vertex uniformly by one of \(r\) labels.
A given \(r\)-clique has all labels distinct with probability
\(r!/r^r\). By averaging there is a partition into the label classes
with at least this proportion of the cliques crossing all classes.
They give transversal tuples in the product of the classes, so
Lemma~\ref{p205:transversal-density} supplies \(t\) points in each
class all of whose transversal tuples are cliques. Any two points
from different chosen classes can be completed to a transversal tuple,
and therefore are adjacent. They give the required subgraph.
\end{proof}

\begin{theorem}[Erd\H{o}s--Stone--Simonovits]\label{p205:W04}
Let \(H\) be a finite simple graph with at least one edge, and
write \(\operatorname{ex}(n,H)\) for the largest number of edges
in an \(n\)-vertex \(H\)-free graph. Then
\[
 \lim_{n\to\infty}\frac{\operatorname{ex}(n,H)}{\binom n2}
   =1-\frac1{\chi(H)-1}.
\]
In particular, the limiting density for bipartite \(H\) is zero.
\end{theorem}
\begin{proof}
Set \(r=\chi(H)\ge2\). The graph \(T_{r-1}(n)\) is \((r-1)\)-colorable,
so it contains no copy of \(H\). If its part sizes are \(a_i\), then
\(a_i/n\to1/(r-1)\) for every \(i\). Its edge count is
\((n^2-\sum_i a_i^2)/2\), giving the required limit inferior.
In a proper \(r\)-coloring of \(H\), every color class has at most
\(t=|V(H)|\) vertices. Map each color class injectively into its
part of \(K_r(t)\) to embed \(H\). For each fixed \(\varepsilon>0\),
Theorem~\ref{p205:W03} then rules out the threshold density
\(1-1/(r-1)+\varepsilon\) for all sufficiently large \(H\)-free
graphs. Letting \(\varepsilon\) tend to zero gives the limit superior.
\end{proof}

The clique bounds and the general density theorem are treated in
\cite[Section~7.1]{diestel} and \cite[Sections~1.3 and~1.5]{zhao}.
The bipartite bound is the classical estimate of \cite{kst}.
\chapter{Extremal Set Families}

For a finite set \(X\), write \(\binom Xk\) for its layer of
\(k\)-element subsets. A family of subsets is an antichain if
distinct members are incomparable by inclusion. It is intersecting
if any two members have a point in common. For a family
\(\mathcal F\subseteq\binom Xk\), its lower shadow \(\partial\mathcal F\)
consists of the \((k-1)\)-sets contained in members of \(\mathcal F\)
when \(k\ge1\). The shadow of a family of empty sets is defined
to be empty.

\section{Antichains and layer densities}

\begin{theorem}[LYM and Sperner]\label{p205:W05}
If \(|X|=n\) and \(\mathcal A\) is an antichain of subsets of \(X\), then
\[
 \sum_{A\in\mathcal A}\binom n{|A|}^{-1}\le1,
 \qquad |\mathcal A|\le\binom n{\lfloor n/2\rfloor}.
\]
Equality in the first inequality holds exactly for full layers.
Equality in the second holds exactly for full middle layers
\(\binom X{\lfloor n/2\rfloor}\) or \(\binom X{\lceil n/2\rceil}\).
\end{theorem}
\begin{proof}
A permutation of \(X\) determines the chain of all its initial
sets, including the empty set and \(X\). A \(k\)-set occurs in
exactly \(k!(n-k)!\) of the \(n!\) permutation chains, counting
permutations with multiplicity. Each chain contains at most one
member of \(\mathcal A\). Counting the pairs of a permutation and
a member on its chain gives the first inequality after division by
\(n!\).

If equality holds, each permutation chain contains exactly one
member. Swapping two adjacent points in a permutation changes
exactly one initial set. If the chosen member changes, its rank
must be that of the changed set. In fact, if the old chosen member
was unchanged, it would still be the unique member on the new
chain. If it was changed, all unchanged sets fail to belong to
the family, so the new chosen member is at the same rank. Any
permutation can be obtained from any other by adjacent swaps,
for example by bringing the desired first point to the front
and repeating on the tail. Hence there is one rank for all
chosen members. Every set at this rank occurs on a permutation
chain and must be chosen there. Conversely each full layer
makes the first inequality an equality.

The ratio of successive binomial coefficients is \((n-k)/(k+1)\)
for \(0\le k<n\), so their maximum is
\(M=\binom n{\lfloor n/2\rfloor}\), attained precisely at the middle
rank or ranks. Each summand is at least \(1/M\). For equality in the
cardinality bound, the sum must equal one and each member must have
a maximizing rank, so the first equality classification gives the
second one. For \(n=0\) there is only the empty-set layer, and for
\(n=1\) both ranks are middle ranks. The count and its equality
conditions apply in both cases.
\end{proof}

\begin{proposition}[Local LYM]\label{p205:W06}
If \(1\le k\le n=|X|\) and \(\mathcal A\subseteq\binom Xk\), then
\[
 \frac{|\partial\mathcal A|}{\binom n{k-1}}
 \ge \frac{|\mathcal A|}{\binom nk}.
\]
If \(0\le k<n\), the upper shadow \(\partial^+\mathcal A\) consisting
of \((k+1)\)-sets containing a member similarly satisfies
\[
 \frac{|\partial^+\mathcal A|}{\binom n{k+1}}
 \ge \frac{|\mathcal A|}{\binom nk}.
\]
\end{proposition}
\begin{proof}
Count pairs \((A,B)\) with \(A\in\mathcal A\), \(B\subset A\) and
\(|B|=k-1\). The count is \(k|\mathcal A|\), and at most
\(n-k+1\) extensions are possible for each \(B\). Thus
\(k|\mathcal A|\le(n-k+1)|\partial\mathcal A|\), which is the first
claim by the ratio of binomial coefficients. Complementation
interchanges the upper shadow of \(\mathcal A\) with the lower
shadow of its complementary family in rank \(n-k\).
\end{proof}

\section{Uniform intersecting families}

\begin{theorem}[Erd\H{o}s--Ko--Rado bound]\label{p205:W07}
For \(1\le k\le n/2\), an intersecting family
\(\mathcal F\subseteq\binom Xk\) satisfies
\[
 |\mathcal F|\le\binom{n-1}{k-1}.
\]
The full star, consisting of all \(k\)-sets containing a fixed
point, attains the bound. If \(n=2k\), equality holds exactly when
the family contains one member of every complementary pair.
\end{theorem}
\begin{proof}
Place the points in an arbitrary cyclic order with positions
numbered modulo \(n\). Consider the arcs of \(k\) consecutive points
that belong to \(\mathcal F\). If at least one belongs, rotate
positions so that one selected arc begins at \(0\). An arc that
intersects this arc has a start in
\(\{-(k-1),\ldots,k-1\}\) modulo \(n\). These \(2k-1\) positions
are distinct because \(n\ge2k\). For each \(i=1,\ldots,k-1\),
the arcs starting at \(i-k\) and \(i\) have disjoint consecutive
blocks, whose combined length is \(2k\le n\). Thus at most one
from each pair is selected, and there are at most \(k\) selected arcs.
If no arc belongs to the family the same bound holds.

Now choose the permutation assigning points to positions uniformly.
For each position its \(k\)-arc is a uniform \(k\)-set. Averaging the
number of selected arcs gives
\(n|\mathcal F|/\binom nk\le k\). The bound follows, and the full
star has precisely that many members by counting its remaining
\(k-1\) points. When \(n=2k\), the only disjoint \(k\)-sets are
complements of one another. At most one can be taken from each
pair, and \(\binom nk/2=\binom{n-1}{k-1}\), proving the equality
criterion.
\end{proof}

\begin{theorem}[Uniqueness in the strict range]\label{p205:W08}
If \(k\ge1\) and \(n>2k\), the full stars are the only intersecting
families of \(k\)-sets of size \(\binom{n-1}{k-1}\).
\end{theorem}
\begin{proof}
Order \(X=\{1,\ldots,n\}\). For \(i<j\), shift a set containing
\(j\) and avoiding \(i\) to the set with \(j\) replaced by \(i\)
provided the replacement is absent from the family. Keep every
other set. This preserves the size since each moved set has a
distinct, previously absent image. It also preserves intersection.
Indeed, if two resulting sets are disjoint, at least one was
moved, and they cannot both have moved because both images
would contain \(i\). Write the moved set as \(C\cup\{i\}\) with
preimage \(C\cup\{j\}\), and the retained set as \(B\). The old
intersection then forces \(j\in B\). Since \(B\) avoids \(i\) and
was retained, \(B\setminus\{j\}\cup\{i\}\) was in the old family,
disjoint from \(C\cup\{j\}\). This is a contradiction.
Perform any shift that changes the family until no such shift remains.
The sum of all point labels over all members strictly decreases at
each change, so this process terminates.

We prove the classification by induction on \(k\), uniformly in
\(n>2k\). At \(k=1\) the maximizing family is a singleton star.
For \(k\ge2\), let \(\mathcal F_0\) be the sets of the shifted
family avoiding \(n\), and let \(\mathcal F_1\) consist of the
sets containing \(n\) with \(n\) removed. Both are intersecting
families on \(n-1\) points. For \(\mathcal F_1\), if two members
\(C,D\) were disjoint, choose \(i\in\{1,\dots,n-1\}\setminus(C\cup D)\).
The shifted family would contain \(C\cup\{i\}\), which is disjoint
from \(D\cup\{n\}\). By Theorem~\ref{p205:W07},
\[
 |\mathcal F_0|\le\binom{n-2}{k-1},
 \qquad |\mathcal F_1|\le\binom{n-2}{k-2}.
\]
Both ranks are within the theorem's range because \(n-1\ge2k\).
Pascal's identity and the extremal size force equality in both.
Since \(n-1>2(k-1)\), induction shows that \(\mathcal F_1\) is a
full star at some \(c<n\). A member of \(\mathcal F_0\) avoiding
\(c\) would be disjoint from a \((k-1)\)-set containing \(c\)
on the remaining points of \(\{1,\ldots,n-1\}\). There are at
least \(k-1\) remaining points, so this is impossible.
Thus the entire shifted family lies in the star at \(c\)
and fills it by its size.

It remains to undo each shift. Suppose the \((i,j)\)-shift of an
intersecting family is a full star at \(c\). If \(c\notin\{i,j\}\),
a shift preserves whether each set contains \(c\), so the original
family is that same star by cardinality. If \(c=j\), take any
\((k-1)\)-set \(D\) avoiding \(i,j\). The set \(D\cup\{j\}\)
cannot survive a shift without \(D\cup\{i\}\) also being present,
which is impossible in the output star. If \(c=i\), the original
family contains all \(k\)-sets containing both \(i,j\), none
containing neither, and exactly one from
\(\{D\cup\{i\},D\cup\{j\}\}\) for each such \(D\). This follows by
checking the shift on each two-set pair, where the number of
members is preserved. For disjoint \(D,D'\) the chosen points
from \(\{i,j\}\) must agree, or the two corresponding members
would be disjoint.

The graph of \((k-1)\)-sets of \(X\setminus\{i,j\}\) joined when
disjoint is connected. To see this, one can connect any two
\((k-1)\)-sets by successive single-point replacements. The union
of two successive sets has \(k\) elements. Its complement in
\(X\setminus\{i,j\}\) has size \(n-2-k\ge k-1\), so these two
sets have a common disjoint neighbor. This gives a walk in the
disjointness graph. The choice between \(i,j\) is therefore
constant and the original family is a full star. Undoing the
finitely many shifts completes the proof.
\end{proof}

\section{The minimum shadow}

Give \(X\) a linear order. For distinct sets \(A,B\) of the same
size, let \(A<_{\mathrm{colex}}B\) if the largest point of
\(A\mathbin{\triangle} B\) belongs to \(B\). This is a linear order
on each layer. In the order \(x_1<\cdots<x_n\) on \(X\), it agrees
with increasing values of \(\sum_{x_i\in A}2^{i-1}\), since the
largest differing term outweighs the sum of the smaller terms.
Let \(I_k(m)\) be the first \(m\) sets in this order.

\begin{lemma}[The shadow of an initial segment]\label{p205:colex-shadow}
For \(1\le k\le n\), the shadow of a colex initial segment in
\(\binom Xk\) is a colex initial segment in \(\binom X{k-1}\).
If the initial segment is nonempty with greatest member \(S\),
then its shadow ends at \(S\setminus\{\min S\}\).
\end{lemma}
\begin{proof}
The empty case is immediate. Put \(S'=S\setminus\{\min S\}\).
For a \(k\)-set \(A\le_{\mathrm{colex}} S\), the greatest of its
\((k-1)\)-subsets is obtained by removing \(\min A\). We show that
this set is at most \(S'\). If \(A=S\), this is clear. Otherwise
let \(d\in S\setminus A\) be the largest differing point.
If \(d=\min S\), the two sets agree above \(d\), with \(A\) having
one more point below \(d\). Removing the two minima therefore
gives the same set. If \(d>\min S\), then \(\min A<d\). Otherwise
\(A\) would have all its \(k\) points above \(d\), where it agrees
with \(S\), although \(S\) also contains \(d\). Removing both minima
then leaves \(d\) as the largest differing point and it lies in \(S'\).
It follows that every point of the shadow is at most \(S'\).

Conversely, let \(B\le_{\mathrm{colex}}S'\) have size \(k-1\).
For \(B=S'\) use the extension \(S\). Otherwise let \(d\in S'\setminus B\)
be the largest differing point, and let \(x\) be the least point of
\(X\setminus B\). Thus \(x\le d\). If \(x<d\), then \(B\cup\{x\}\)
is below \(S\) in colex, since \(\min S<d\) and \(d\) is still the
largest differing point. If \(x=d\), then \(B\) contains every point
below \(d\). Above \(d\) the sets \(B,S'\) agree. Since \(S'\) has
the same size as \(B\) and contains \(d\), it has one fewer point
below \(d\). The additional point \(\min S\) fills this single
gap below \(d\), and \(S=B\cup\{d\}\). Thus every such \(B\)
belongs to the shadow.
\end{proof}

For disjoint, equally sized sets \(U,V\subseteq X\), call \(A\)
movable when it contains \(V\) and avoids \(U\). Its replacement
is \(A^*=(A\setminus V)\cup U\). The compression \(C_{U,V}\)
of a family replaces each movable member whose replacement
is absent, and leaves other members in place. The inverse on
replacements removes \(U\) and restores \(V\). In particular the
compression preserves rank and cardinality. A family is
\(C_{U,V}\)-closed if it contains the replacement of each
movable member. For \(U=V=\varnothing\) this holds automatically.

\begin{lemma}[Smaller compressions control the shadow]\label{p205:uv-shadow}
Suppose \(U,V\) are disjoint nonempty sets of equal size. For every
\(x\in U\), assume that some \(y\in V\) has the property that
\(\mathcal F\) is \(C_{U\setminus\{x\},\,V\setminus\{y\}}\)-closed.
Then
\[
 \partial(C_{U,V}\mathcal F)
   \ \subseteq\ C_{U,V}(\partial\mathcal F).
\]
In particular, \(C_{U,V}\) does not increase the shadow size.
\end{lemma}
\begin{proof}
Take \(B\) in the left side and write \(A=B\cup\{z\}\) for an
extension in \(C_{U,V}\mathcal F\), with \(z\notin B\).
First suppose \(A\notin\mathcal F\). Then
\(A=D\cup U\) has been moved from \(D\cup V\in\mathcal F\),
where \(D\) avoids \(U\cup V\). If \(z\notin U\),
\((D\setminus\{z\})\cup V\) belongs to \(\partial\mathcal F\)
and its replacement is \(B\). Hence \(B\) belongs to the
compressed shadow, whether it is newly moved there or already
present. If \(z\in U\), choose \(y\in V\) as in the hypothesis
for \(x=z\). Compressing \(D\cup V\) along
\((U\setminus\{z\}, V\setminus\{y\})\) gives
\[
   D\cup(U\setminus\{z\})\cup\{y\}=B\cup\{y\}\in\mathcal F.
\]
Thus \(B\in\partial\mathcal F\) and it is retained under \(C_{U,V}\)
because it does not contain \(V\).

Now suppose \(A\in\mathcal F\). Then \(B\in\partial\mathcal F\), and
we need consider only movable \(B\), say \(B=D\cup V\)
with \(D\) avoiding \(U\cup V\). If \(z\notin U\),
\(A\) is also movable. It survived in the compressed family, so its
replacement \(D\cup U\cup\{z\}\) was in \(\mathcal F\).
The set \(B^*=D\cup U\) therefore belongs to \(\partial\mathcal F\).
If \(z\in U\), choose \(y\in V\) for \(x=z\). The smaller
compression sends \(D\cup V\cup\{z\}\) to
\(D\cup U\cup\{y\}\), which lies in \(\mathcal F\) by closure.
Again \(B^*\in\partial\mathcal F\). In both subcases \(B\) is
retained because its replacement was already in the shadow.
The inclusion follows. The final inequality uses cardinality
preservation of the compression applied to the shadow.
\end{proof}

\begin{theorem}[Kruskal--Katona in colex form]\label{p205:W09}
Let \(1\le k\le |X|\), let \(\mathcal F\subseteq\binom Xk\), and
put \(m=|\mathcal F|\). For every \(0\le j\le k\),
\[
   \bigl|\partial^j I_k(m)\bigr|\le\bigl|\partial^j\mathcal F\bigr|,
\]
where \(\partial^0\) is the identity.
\end{theorem}
\begin{proof}
Call a pair of nonempty disjoint equally sized sets \(U,V\) useful
if \(\max U<\max V\). If some useful compression changes
\(\mathcal F\), choose such a pair with \(|U|\) minimal. If
\(|U|=1\), deleting the points gives the identity compression.
Otherwise let \(y=\min V\). For every \(x\in U\) the pair
\(U\setminus\{x\}, V\setminus\{y\}\) is still useful since
\(\max(V\setminus\{y\})=\max V\). Its compression fixes
\(\mathcal F\) by minimality. A compression fixes the family
exactly when it is closed. Thus Lemma~\ref{p205:uv-shadow}
shows that the chosen compression does not increase the shadow.

Each changed member decreases in colex because the largest point
of its symmetric difference with its replacement belongs to \(V\).
The sum of the colex ranks of the family members is a nonnegative
integer, strictly decreasing at each step. Repeating the choice
eventually gives a family closed under every useful compression.
Such a family is a colex initial segment. Indeed, if \(A\) belongs
and \(B<_{\mathrm{colex}}A\) is absent, then
\(U=B\setminus A,\ V=A\setminus B\) form a useful pair sending
\(A\) to \(B\), contradicting closure. Rank and cardinality
were preserved, so the final family is \(I_k(m)\) and
\(|\partial I_k(m)|\le|\partial\mathcal F|\).

Now suppose \(\partial^j I_k(m)\) is an initial segment of size
\(a\le b=|\partial^j\mathcal F|\) in rank \(k-j>0\).
It is contained in the initial segment of size \(b\) in that rank.
Taking shadows preserves inclusion, and the one-step inequality
for \(\partial^j\mathcal F\) bounds the shadow of the larger
initial segment. Lemma~\ref{p205:colex-shadow} ensures that
\(\partial^{j+1} I_k(m)\) is again an initial segment. Induction
starting at \(j=0\) proves the assertion through \(j=k\).
\end{proof}

\begin{remark}
The argument uses an arbitrary finite linearly ordered ground set.
Its increasing enumeration identifies it with the ordered \(n\)-point
set used by the Mathlib colex declarations. An order-preserving
bijection sends the greatest point of any symmetric difference
to the greatest point of its image, preserves ranks, and sends
one-point deletions bijectively to one-point deletions.
Thus it transports both colex segments and every iterated shadow.
The colex compression route follows the structure of the checked
Mathlib proof~\cite{mathlib}. The binomial expansion formulation
is discussed in \cite[Section~10.4]{jukna}.
\end{remark}

\section{Set-pair inequalities}

\begin{theorem}[Bollob\'as's set-pair inequality]\label{p205:W10}
Let \((A_i,B_i)\) be finitely many pairs of finite sets. Assume
\(A_i\cap B_i=\varnothing\) for every \(i\) and
\(A_i\cap B_j\ne\varnothing\) for all ordered pairs \(i\ne j\).
Then
\[
   \sum_i\binom{|A_i|+|B_i|}{|A_i|}^{-1}\le1.
\]
\end{theorem}
\begin{proof}
Choose a uniform permutation of the union of all sets. Let \(E_i\)
be the event that every point of \(A_i\) precedes every point of
\(B_i\). The relative order of \(A_i\cup B_i\) is uniform.
Choosing the positions of \(A_i\) gives
\(\mathbb P(E_i)=\binom{|A_i|+|B_i|}{|A_i|}^{-1}\), including empty
sides. If \(i\ne j\), choose \(x\in A_i\cap B_j\) and
\(y\in A_j\cap B_i\). These points are distinct by the within-pair
disjointness. Event \(E_i\) requires \(x\) before \(y\), and \(E_j\)
requires \(y\) before \(x\), so the events are pairwise disjoint.
Adding their probabilities proves the claim.
\end{proof}

For the classical antichain, intersecting-family, and set-pair
bounds, see \cite[Sections~7.2 and~8.3--8.4]{jukna}.
\chapter{Finite and Infinite Ramsey Theory}

For a set \(X\) and a positive integer \(k\), write \([X]^k\) for
its \(k\)-element subsets. A subset \(Y\subseteq X\) is homogeneous
for a coloring of \([X]^k\) if there is a color that all members of
\([Y]^k\) have. For a nonempty color set this condition holds
vacuously whenever \(|Y|<k\).

\section{Infinite and finite multicolor Ramsey}

\begin{theorem}[Infinite Ramsey]\label{p205:W11}
Let \(k,r\ge1\) be integers. Every coloring of \([\mathbb N]^k\)
with \(r\) colors has an infinite homogeneous subset. The same
conclusion holds for a coloring of the \(k\)-sets of any infinite set.
\end{theorem}
\begin{proof}
For \(k=1\), an infinite set partitioned into finitely many
parts has an infinite part, since a finite union of finite sets is
finite. Suppose the result holds for \(k-1\), where \(k\ge2\).
The induction hypothesis also holds on each infinite subset of
\(\mathbb N\) by transport along its increasing enumeration.

Starting with \(S_0=\mathbb N\), recursively take \(a_i=\min S_i\).
The points of \(S_i\) greater than \(a_i\) form an infinite set.
Color their \((k-1)\)-sets \(U\) by the old color of
\(\{a_i\}\cup U\). By induction this induced coloring is constant,
with some color \(d_i\), on the \((k-1)\)-sets of an infinite
subset \(S_{i+1}\subseteq S_i\setminus\{a_i\}\). The \(a_i\) are
strictly increasing. Among the indices \(i\) there is an infinite
set \(I\) with a common \(d_i=d\). For any \(k\)-set drawn from
\(\{a_i:i\in I\}\), let \(a_i\) be its least point. All the other
points belong to \(S_{i+1}\) by nesting, so the \(k\)-set has color \(d\).
This proves homogeneity. For an arbitrary infinite set, choose a
sequence of distinct points recursively from the complements of its
finite initial segments. Pull back along this injection of
\(\mathbb N\), and take the image of an infinite homogeneous subset.
\end{proof}

\begin{theorem}[Finite Ramsey]\label{p205:W12}
For integers \(k,r,m\ge1\) there is an integer \(N\ge m\) such that
every \(r\)-coloring of the \(k\)-sets of an \(N\)-point set has a
homogeneous \(m\)-set.
\end{theorem}
\begin{proof}
Take \(N=m\) if \(m<k\) or \(r=1\). For \(k=1\), take
\(N=r(m-1)+1\) by counting points in the color classes.
In the remaining cases suppose that no \(N\) works. Call a coloring
of the \(k\)-sets of \(\{1,\ldots,n\}\) bad if it has no homogeneous
\(m\)-set. Every restriction of a bad coloring to a shorter
initial segment is bad. Organize the bad colorings into a tree by
connecting each level \(n+1\) coloring to its restriction at level
\(n\). The level-zero root has bad extensions of arbitrarily large
length, and every node has finitely many children since a finite
set of new \(k\)-sets must be colored with finitely many colors.

At a node with arbitrarily long extensions, at least one child
has arbitrarily long extensions. Otherwise take the maximum of
the finitely many upper bounds for its children to obtain a bound
at the node. Repeatedly choose such a child to obtain compatible
colorings at all levels. For each \(k\)-set of positive integers,
its color is fixed by every sufficiently long initial segment,
so these colorings define an infinite coloring. An infinite
homogeneous subset exists by Theorem~\ref{p205:W11}. Any \(m\)
of its points lie in some finite initial segment, contradicting
that segment's bad coloring.
\end{proof}

\section{Two-color graph Ramsey numbers}

\begin{theorem}[Recursion and the first diagonal value]\label{p205:W13}
For integers \(a,b\ge1\), let \(R(a,b)\) be the least positive \(n\)
such that every red--blue edge coloring of \(K_n\) has a red \(K_a\)
or a blue \(K_b\). Then \(R(1,b)=R(a,1)=1\), and for \(a,b\ge2\),
\[
 R(a,b)\le R(a-1,b)+R(a,b-1)
    \qquad\text{and}\qquad
 R(a,b)\le\binom{a+b-2}{a-1}.
\]
In particular \(R(3,3)=6\).
\end{theorem}
\begin{proof}
Existence follows from Theorem~\ref{p205:W12} with \(k=r=2\)
and \(m=\max\{a,b\}\). A single vertex is a clique in either
color, proving the boundary values. Put \(u=R(a-1,b)\),
\(v=R(a,b-1)\), and fix a vertex \(x\) in a colored \(K_{u+v}\).
Among the \(u+v-1\) other vertices, either at least \(u\) are
joined to \(x\) in red, or at least \(v\) in blue. In the red case
the \(u\) chosen neighbors contain a blue \(K_b\) or a red
\(K_{a-1}\), which extends by \(x\). In the blue case the symmetric
argument yields a red \(K_a\) or a blue \(K_b\).
The binomial bound now follows by induction on \(a+b\) using
Pascal's identity with the boundary values. It gives
\(R(3,3)\le\binom42=6\). On five cyclically ordered vertices, color
the cycle edges red and the other edges blue. Each color forms a
five-cycle, so neither contains a triangle. Thus \(R(3,3)>5\).
\end{proof}

The general Ramsey statements follow the distinction between infinite
induction and finite compactness in \cite[Section~9.1]{diestel}.
For the two-color graph bounds see \cite[Section~4.9]{jukna}.
\chapter[Elementary Probabilistic Method]{The Elementary Probabilistic Method}
\label{p205:probabilistic}

A probability argument can find a combinatorial object by studying the
distribution of a sample that encodes it. The witness is an outcome whose
decoded object has the desired property. We develop avoidance, alteration,
and a deterministic selection rule on finite experiments. Alon and Spencer
give the classical starting points for these methods \cite{alon}.

\section{Positive probability and finite avoidance}

Write \((\Omega,\mathcal F,\mathbb P)\) for a probability space. In a finite
weighted experiment we take \(\mathcal F=2^\Omega\) and give each outcome
\(\omega\) a weight \(p_\omega\ge 0\), with
\(\sum_{\omega\in\Omega}p_\omega=1\).
Then \(\mathbb P(A)=\sum_{\omega\in A}p_\omega\). The positive support is
\(\{\omega\in\Omega:p_\omega>0\}\). If \(Z:\Omega\to\mathbb R\) is a real
quantity in such an experiment, its expectation is
\(\mathbb E Z=\sum_{\omega\in\Omega}p_\omega Z(\omega)\).
For an event \(A\), the indicator \(\mathbf 1_A\) equals \(1\) on \(A\) and
\(0\) on its complement, so \(\mathbb E\mathbf 1_A=\mathbb P(A)\). The sum
formula makes expectation linear even if the summed quantities are
dependent.

\begin{proposition}[Existence from positive probability]
\label{p205:W14}
Let \(A\in\mathcal F\) satisfy \(\mathbb P(A)>0\). Then \(A\) is nonempty.
In a finite weighted experiment, \(A\) contains a positive-weight
outcome. In particular, if \(f\) sends outcomes to combinatorial objects
and the event \(A\) consists of outcomes whose images under \(f\) have a
specified property, that property is realized by an object \(f(\omega)\).
\end{proposition}

\begin{proof}
The empty set has probability zero, so \(A\) contains an outcome.
In the finite case, a positive sum of nonnegative weights has a positive
summand. Apply \(f\) to the chosen outcome.
\end{proof}

\begin{proposition}[Union bound and simultaneous avoidance]
\label{p205:W15}
For measurable events \(A_i\), indexed by a finite set \(I\), one has
\[
  \mathbb P\!\left(\bigcup_{i\in I} A_i\right)
       \le \sum_{i\in I}\mathbb P(A_i).
\]
If the sum on the right is strictly less than \(1\), then
\(\mathbb P(\bigcap_{i\in I} A_i^c)>0\).
\end{proposition}

\begin{proof}
Enumerate \(I\) as \(i_1,\ldots,i_m\) and define
\[
 B_j=A_{i_j}\setminus\bigcup_{\ell<j}A_{i_\ell}.
\]
These measurable sets are pairwise disjoint, have the same union as the
\(A_i\), and satisfy \(B_j\subseteq A_{i_j}\).
Finite additivity and monotonicity give
the inequality. Taking complements shows that the probability of
simultaneous avoidance is at least \(1-\sum_i\mathbb P(A_i)>0\).
For \(I=\varnothing\) the bad union is empty and the good intersection
is \(\Omega\), with probability \(1\).
\end{proof}

\begin{theorem}[A random-coloring lower bound for Ramsey numbers]
\label{p205:W16}
For integers \(n\ge 0\) and \(k\ge 2\), if
\[
  \binom nk\,2^{\,1-\binom k2}<1,
\]
then \(R(k,k)>n\). In particular \(R(k,k)>2^{k/2}\) for \(k\ge3\).
\end{theorem}

\begin{proof}
Color each edge of \(K_n\) red or blue independently with probability
\(1/2\) for either color. For any fixed \(k\)-set of vertices, there are
\(\binom k2\) edges on it and the two monochromatic patterns are
distinct, since \(k\ge2\). Its bad-event probability is
\(2^{1-\binom k2}\). By Proposition~\ref{p205:W15}, a coloring exists
that has no monochromatic \(K_k\). If the Ramsey number were at most
\(n\), restricting that coloring to \(R(k,k)\) vertices would give a
contradiction. When \(n<k\), the family of bad events is empty and the
same conclusion holds.

Take \(n=\lfloor 2^{k/2}\rfloor\). Since \(\binom nk\le n^k/k!\), the
left side of the criterion is at most \(2^{1+k/2}/k!\). This is strictly
less than \(1\) at \(k=3\), since \(2^{5/2}<6\). Increasing \(k\) to
\(k+1\) multiplies the bound by \(\sqrt{2}/(k+1)<1\). Therefore
\(R(k,k)\ge \lfloor 2^{k/2}\rfloor+1>2^{k/2}\).
\end{proof}

\section{Alteration and deterministic choices}

\begin{theorem}[The finite alteration principle]
\label{p205:W17}
For each outcome of a finite weighted experiment, suppose a finite
object has size \(S\) and an associated nonnegative integer defect
count \(D\). Suppose a specified deterministic deletion procedure
repairs every object on the positive support, removing at most \(D\)
items, and yields an object of size at least \(S-D\) with all defects
cleared. Then some repaired object has size at least
\(\mathbb E(S-D)\).
\end{theorem}

\begin{proof}
The maximum of \(S-D\) over the nonempty positive support is at least
its weighted mean. Choose an outcome attaining the maximum and apply
the repair procedure. In an application where defects are forbidden
subsets and deletion cannot create new forbidden subsets, a concrete
procedure is to enumerate the original defects and, whenever the
next still survives, delete a member of it. At most one deletion is
charged to each original defect, and every remaining defect would
have been among the original ones and would have been hit at its
turn. An application with different defect data must establish its
own repair property.
\end{proof}

\begin{corollary}[An independence bound by deletion]
\label{p205:W18}
Let \(G\) be a finite simple graph with \(n\) vertices and \(m\) edges.
For each real \(p\in[0,1]\), its independence number satisfies
\[
  \alpha(G)\ge pn-p^2m .
\]
\end{corollary}

\begin{proof}
Independently keep each vertex with probability \(p\). Let \(S\) be
the number of kept vertices and \(D\) the number of edges whose two
endpoints are kept. Their expectations are \(pn\) and \(p^2m\), by
the indicator sum and by independence of the two choices on each
edge. Enumerate the edges originally kept and remove one endpoint
whenever the edge still survives. The remaining vertices are
independent and at most \(D\) vertices were deleted. Apply
Theorem~\ref{p205:W17}. The experiment and deletion argument also
apply at \(p=0\) and \(p=1\) and when the graph is empty.
\end{proof}

\begin{theorem}[Method of conditional expectations]
\label{p205:W19}
In a finite weighted experiment revealed in finitely many stages,
there is a positive-weight outcome whose real objective is at
least its initial expectation. One may choose such an outcome
by selecting a positive-probability branch of nondecreasing
conditional expectation at each stage.
In particular a finite simple graph with \(m\) edges admits a
cut of at least \(m/2\) edges, obtained by this rule on successive
vertex choices.
\end{theorem}

\begin{proof}
A finite staged experiment can be described by a sequence of
successively refining partitions of its outcomes, starting at the
trivial partition and ending at singletons. At a current cell \(C\)
of positive probability define the conditional expectation to be
\[
  \mathbb E[Z\mid C]=
     \frac{\sum_{\omega\in C}p_\omega Z(\omega)}{\mathbb P(C)}.
\]
If \(C_1,\ldots,C_s\) are its positive-probability child cells, their
total probability is \(\mathbb P(C)\). The defining sum split among
the children yields
\[
  \mathbb E[Z\mid C]
     =\sum_{j=1}^{s}\frac{\mathbb P(C_j)}{\mathbb P(C)}
                      \mathbb E[Z\mid C_j].
\]
At least one child's conditional expectation is at least the
parent's. Choose such a child at every stage. At the terminal
singleton the conditional expectation is the objective value.

For the cut, color all vertices with independent fair bits and let
\(Z\) count edges whose ends receive different bits. Each edge is cut
with probability \(1/2\), so \(\mathbb E Z=m/2\). Reveal one bit at a
time. On any prefix the contribution of an edge with both ends
revealed is its known indicator. If at least one endpoint is
unrevealed, the contribution is \(1/2\), since the unrevealed
choices remain independent and fair. This supplies the conditional
means used to choose successive bits. An empty graph has a cut with
zero edges and the same argument permits an empty sequence of choices.
\end{proof}
\chapter{Moments and Concentration}
\label{p205:moments}

This chapter moves from expected counts to quantitative probability
bounds. The moment inequalities are stated on probability spaces, so
that the same formulas apply to finite experiments and to measurable
random variables with the indicated integrability. The finite
applications require only weighted sums.

\section{Indicator variables, first moments, and second moments}

For a real random variable \(X\) we write \(\mathbb E X\) for its
integral under the probability measure. If \(X\) and \(Y\) have finite
second moments, put
\[
 \operatorname{Var}(X)=\mathbb E[(X-\mathbb E X)^2],\qquad
 \operatorname{Cov}(X,Y)=
 \mathbb E[(X-\mathbb E X)(Y-\mathbb E Y)] .
\]
The expression for the covariance is integrable, since
\(2|uv|\le u^2+v^2\) for real \(u,v\). Expanding the square and the
product gives
\(\operatorname{Var}(X)=\mathbb E[X^2]-(\mathbb E X)^2\) and
\(\operatorname{Cov}(X,Y)=\mathbb E[XY]-\mathbb E X\,\mathbb E Y\).

\begin{proposition}[Indicator counts and the Caro--Wei bound]
\label{p205:W20}
Let \(A_i\), \(i\in I\), be a finite family of measurable events, and
let \(Z=\sum_i\mathbf 1_{A_i}\), with \(p_i=\mathbb P(A_i)\). Then
\[
 \mathbb E Z=\sum_i p_i,\qquad
 \operatorname{Var}(Z)=
  \sum_{i,j\in I}\operatorname{Cov}(\mathbf 1_{A_i},\mathbf 1_{A_j}).
\]
In an enumeration of \(I\), the second formula is
\[
 \operatorname{Var}(Z)=\sum_i p_i(1-p_i)
   +2\sum_{i<j}\bigl(\mathbb P(A_i\cap A_j)-p_i p_j\bigr).
\]
If \(G\) is a finite simple graph and \(d(v)\) is the degree of \(v\),
then
\[
 \alpha(G)\ge\sum_{v\in V(G)}\frac{1}{d(v)+1}.
\]
\end{proposition}

\begin{proof}
The expectation formula follows by linearity. Write
\(Z-\mathbb E Z=\sum_i(\mathbf 1_{A_i}-p_i)\), expand its square and
integrate. For the diagonal terms use \(\mathbf 1_{A_i}^2=\mathbf
1_{A_i}\). For each off-diagonal term use
\(\mathbf 1_{A_i}\mathbf 1_{A_j}=\mathbf 1_{A_i\cap A_j}\).
In particular the off-diagonal covariance is zero whenever the
two events are independent.

For the graph bound, choose a uniform random ordering of the
vertices. Include each vertex that precedes all its neighbors.
The selected vertices form an independent set, since adjacent
vertices cannot both precede one another. The relative order on
\(\{v\}\cup N(v)\) is uniform: each of its \((d(v)+1)!\) orders
extends to the same number of orders of the full vertex set.
Thus the selection probability of \(v\) is \(1/(d(v)+1)\).
The expected size of the selected set equals the stated sum,
and every selected set has size at most \(\alpha(G)\).
The sum is zero when the vertex set is empty.
\end{proof}

\begin{theorem}[Markov's inequality and the zero-defect criterion]
\label{p205:W21}
If \(X\ge0\) almost surely is integrable and \(a>0\), then
\[
 \mathbb P(X\ge a)\le \frac{\mathbb E X}{a}.
\]
If \(X\) is nonnegative integer-valued almost surely and
\(\mathbb E X<1\), then \(\mathbb P(X=0)>0\).
\end{theorem}

\begin{proof}
The pointwise inequality \(a\mathbf 1_{\{X\ge a\}}\le X\) holds
almost surely. Integrating and dividing by \(a>0\) gives the
first assertion. Under the integer-valued hypothesis,
\(\{X\ne0\}\) and \(\{X\ge1\}\) agree off a null set. The
first assertion with \(a=1\) bounds their probability strictly
below \(1\).
\end{proof}

\begin{theorem}[Second-moment existence bound]
\label{p205:W22}
If \(X\ge0\) almost surely, \(0<\mathbb E X<\infty\), and
\(\mathbb E[X^2]<\infty\), then
\[
  \mathbb P(X>0)\ge
       \frac{(\mathbb E X)^2}{\mathbb E[X^2]} .
\]
\end{theorem}

\begin{proof}
For a measurable event \(B\) with \(\mathbb P(B)>0\), the
nonnegative square
\(\mathbb E[(X\mathbf 1_B-c\mathbf 1_B)^2]\) with
\(c=\mathbb E[X\mathbf 1_B]/\mathbb P(B)\) gives
\[
  (\mathbb E[X\mathbf 1_B])^2\le
      \mathbb P(B)\,\mathbb E[X^2\mathbf 1_B]
      \le \mathbb P(B)\,\mathbb E[X^2].
\]
The same inequality holds for \(\mathbb P(B)=0\) because the
integral over a null set vanishes. Use \(B=\{X>0\}\). Its
truncated expectation is \(\mathbb E X>0\), so the inequality
also gives \(\mathbb E[X^2]>0\) before division.
This is the second-moment argument in its indicator form
\cite{roch}.
\end{proof}

\begin{theorem}[Strict Paley--Zygmund and the variance method]
\label{p205:W23}
Let \(X\) satisfy the hypotheses of Theorem~\ref{p205:W22} and
write \(\mu=\mathbb E X\). For \(0\le\theta<1\),
\[
 \mathbb P(X>\theta\mu)\ge
   \frac{(1-\theta)^2\mu^2}{\mathbb E[X^2]}.
\]
Also
\[
 \mathbb P(X=0)\le\frac{\operatorname{Var}(X)}{\mu^2}.
\]
Consequently, for any sequence of such variables, on possibly
different probability spaces, if
\(\operatorname{Var}(X_n)/(\mathbb E X_n)^2\to0\), then
\(\mathbb P(X_n>0)\to1\).
\end{theorem}

\begin{proof}
Let \(B=\{X>\theta\mu\}\). On \(B^c\) the bound is \(X\le\theta\mu\),
so \(\mathbb E[X\mathbf 1_{B^c}]\le\theta\mu\,\mathbb P(B^c)
\le\theta\mu\). Thus \(\mathbb E[X\mathbf 1_B]\ge(1-\theta)\mu\).
Apply the indicator Cauchy--Schwarz inequality from the proof
of Theorem~\ref{p205:W22}, square the positive lower bound, and
divide by \(\mathbb E[X^2]>0\). The strict threshold is important:
the complement used here includes equality. For \(\theta=0\)
one recovers the second-moment bound \cite{tao,alon}.

For any \(t>0\), Markov's inequality applied to
\((X-\mu)^2\) gives
\[
 \mathbb P(|X-\mu|\ge t)
  =\mathbb P((X-\mu)^2\ge t^2)
  \le \frac{\operatorname{Var}(X)}{t^2}.
\]
The event \(X=0\) has deviation exactly \(\mu>0\), and is
included in the deviation event with threshold \(\mu\). This
proves the bound at zero. As \(X_n\ge0\) almost surely,
\(\mathbb P(X_n>0)=1-\mathbb P(X_n=0)\), proving the limit.
\end{proof}

\section{An alteration example}

\begin{theorem}[High girth and high chromatic number]
\label{p205:W24}
Given integers \(g\ge3\) and \(k\ge1\), there is a finite
simple graph with no cycle of length at most \(g\) and with
chromatic number greater than \(k\).
\end{theorem}

\begin{proof}
Fix \(g,k\) and choose \(n\) large. In the random graph on \(n\)
labeled vertices, choose each edge independently with
probability \(p=n^{-1+1/(2g)}\). Let \(C\) count all cycles of
length between \(3\) and \(g\). A cycle of length \(\ell\)
is specified by an ordered list of \(\ell\) distinct vertices,
up to its \(2\ell\) presentations by changing the starting
vertex and direction. Each such cycle occurs with probability
\(p^\ell\). Therefore
\[
 \mathbb E C\le \sum_{\ell=3}^g \frac{n^\ell p^\ell}{2\ell}
      \le g n^{1/2}.
\]
By Theorem~\ref{p205:W21},
\(\mathbb P(C\ge n/2)\le 2g\,n^{-1/2}\).

Put \(a=\lceil n/(4k)\rceil\). For any fixed \(a\) vertices,
independence requires that all \(\binom a2\) possible internal
edges be absent. Proposition~\ref{p205:W15} gives
\[
 \mathbb P(\alpha(G)\ge a)
    \le \binom na(1-p)^{\binom a2}
    \le 2^n\exp\!\bigl(-p\binom a2\bigr).
\]
For \(n\ge 8k\), the inequalities \(a\ge n/(4k)\) and
\(a-1\ge n/(8k)\) give
\[
  p\binom a2\ge \frac{n^{1+1/(2g)}}{64k^2}.
\]
Thus the upper bound for \(\mathbb P(\alpha(G)\ge a)\) tends
to zero as \(n\) tends to infinity. The two bad-event
probabilities thus sum to less than \(1\) for some \(n\).
Choose a graph with \(C<n/2\) and \(\alpha(G)<a\).

Process its original short cycles, deleting a vertex whenever
the next cycle still survives. Deletion creates no new cycle,
and fewer than \(n/2\) vertices are deleted. The remaining
induced graph has more than \(n/2\) vertices and no short
cycle. Every independent set in it was independent in the
original graph, so a color class has size at most
\(a-1<n/(4k)\). At most \(k\) such classes cover fewer than
\(n/4\) vertices. This proves the claimed chromatic number.
The construction is the standard random-graph use of alteration
\cite{alon}.
\end{proof}

\section{Exponential moments}

For a bounded real random variable \(X\), its moment-generating
function is \(M_X(\lambda)=\mathbb E e^{\lambda X}\).
For independent bounded real variables \(X_i\) and bounded
measurable functions \(f_i\) on their essential ranges,
\(\mathbb E\prod_i f_i(X_i)=\prod_i\mathbb E f_i(X_i)\).
For indicator functions this is the defining finite-product
probability identity for independence. Expanding finite sums
extends it to simple functions. Uniform approximation of
bounded measurable functions by simple functions obtained
by partitioning their ranges gives the general identity.
In particular \(M_{\sum_i X_i}(\lambda)=\prod_iM_{X_i}(\lambda)\).
By Markov's inequality, for \(\lambda>0\) and a bounded
real \(Y\),
\[
 \mathbb P(Y\ge t)\le e^{-\lambda t}\mathbb E e^{\lambda Y}.
\]
This follows by applying Theorem~\ref{p205:W21} to the
nonnegative variable \(e^{\lambda Y}\) at \(e^{\lambda t}>0\).

\begin{theorem}[Hoeffding's bound for independent bounded variables]
\label{p205:W25}
Let \(X_1,\ldots,X_m\) be independent real variables with
\(a_i\le X_i\le b_i\) almost surely, for finite real endpoints
\(a_i\le b_i\). Define
\[
  S=\sum_{i=1}^m(X_i-\mathbb E X_i),\qquad
  V=\sum_{i=1}^m(b_i-a_i)^2.
\]
If \(t>0\) and \(V>0\), each of
\(\mathbb P(S\ge t)\) and \(\mathbb P(S\le -t)\) is at most
\(\exp(-2t^2/V)\). If \(V=0\), then \(S=0\) almost surely and
both tail probabilities are zero for \(t>0\).
\end{theorem}

\begin{proof}
First let \(a\le X\le b\) almost surely and set
\(Y=X-\mathbb E X\). Its essential range lies in a bounded
interval \([\alpha,\beta]\) of length \(b-a\).
Let \(M(\lambda)=\mathbb E e^{\lambda Y}\) and
\(f(\lambda)=\log M(\lambda)\). The function \(M\) is strictly
positive. To differentiate, fix a compact interval of
\(\lambda\)-values and an essential bound \(|Y|\le B\).
For \(j=0,1,2\), the absolute value of
\(Y^j e^{\lambda Y}\) is bounded by a finite constant
depending only on \(B\) and the interval. Difference quotients
of the integrands for the first and second derivatives are
bounded by the corresponding derivative bounds on a slightly
larger interval, by the mean-value theorem. Dominated convergence
\cite[Theorem~D11.06.dominated\_convergence]{d11}
yields \(M'(\lambda)=\mathbb E[Ye^{\lambda Y}]\) and
\(M''(\lambda)=\mathbb E[Y^2e^{\lambda Y}]\), with both
derivatives continuous.

The tilted probability measure with density
\(e^{\lambda Y}/M(\lambda)\) has total mass \(1\). The
logarithmic derivative formula gives
\[
 f''(\lambda)=
     \frac{\mathbb E[Y^2e^{\lambda Y}]}{M(\lambda)}
     -\left(\frac{\mathbb E[Ye^{\lambda Y}]}{M(\lambda)}\right)^2,
\]
the variance under this measure. For its expectation
\(\mathbb E_\lambda\), any constant \(u\) satisfies
\[
 \mathbb E_\lambda[(Y-u)^2]
  =\operatorname{Var}_\lambda(Y)+(\mathbb E_\lambda Y-u)^2 .
\]
Choose \(u=(\alpha+\beta)/2\). Almost surely
\((Y-u)^2\le (b-a)^2/4\), and the tilted measure also has this
almost-sure bound. Hence \(f''(\lambda)\le (b-a)^2/4\) for all
real \(\lambda\). Since \(f(0)=f'(0)=0\), integrating twice
along the line segment from \(0\) to \(\lambda\) gives
\[
 f(\lambda)=\lambda^2\int_0^1 (1-s) f''(s\lambda)\,ds
     \le\frac{\lambda^2(b-a)^2}{8},
\]
valid for either sign of \(\lambda\). This proves the
one-variable exponential bound underlying Hoeffding's
inequality \cite{hoeffding,alon}.

Independence now yields
\(\mathbb E e^{\lambda S}\le e^{\lambda^2 V/8}\). The
exponential Markov bound gives
\(\mathbb P(S\ge t)\le\exp(-\lambda t+\lambda^2V/8)\)
for \(\lambda>0\). At \(\lambda=4t/V\) the exponent is
\(-2t^2/V\). Use the same argument for \(-S\) to bound
\(\mathbb P(S\le-t)\). If \(V=0\), all \(a_i=b_i\), so
each centered variable is zero almost surely and their finite
sum is zero almost surely. This also includes the empty sum.
\end{proof}

\pagebreak
\begin{theorem}[Multiplicative Chernoff bounds for Bernoulli sums]
\label{p205:W26}
Let \(X_1,\ldots,X_m\) be independent Bernoulli variables,
with success probabilities \(p_i\in[0,1]\). Put
\(S=\sum_i X_i\) and \(\mu=\sum_i p_i\). For \(\delta\ge0\),
\[
 \mathbb P(S\ge (1+\delta)\mu)
 \le\left(\frac{e^\delta}{(1+\delta)^{1+\delta}}\right)^\mu .
\]
For \(0\le\delta<1\),
\[
 \mathbb P(S\le (1-\delta)\mu)
 \le\left(\frac{e^{-\delta}}{(1-\delta)^{1-\delta}}\right)^\mu .
\]
At \(\delta=1\), the lower-tail bound is
\(\mathbb P(S=0)\le e^{-\mu}\).
If \(\mu=0\), then \(S=0\) almost surely and these
inequalities hold with right sides equal to \(1\).
\end{theorem}

\begin{proof}
For any real \(\lambda\),
\[
 \mathbb E e^{\lambda S}
  =\prod_i(1+p_i(e^\lambda-1))
  \le \exp(\mu(e^\lambda-1)).
\]
Each factor is positive, and the inequality follows from
\(1+u\le e^u\), obtained for example by the minimum of
\(e^u-1-u\) at \(u=0\). For \(\mu>0\) and \(\delta>0\),
use exponential Markov for the upper tail with
\(\lambda=\log(1+\delta)>0\). The exponent is
\(\mu(\delta-(1+\delta)\log(1+\delta))\), which yields the
first formula. If \(0<\delta<1\), take
\(\lambda=\log(1-\delta)<0\) and observe that on the
lower-tail event \(e^{\lambda S}\ge
e^{\lambda(1-\delta)\mu}\). Markov now gives the exponent
\(\mu(-\delta-(1-\delta)\log(1-\delta))\).
For \(\delta=0\) both asserted bounds are \(1\).
For \(\delta=1\), independence gives
\(\mathbb P(S=0)=\prod_i(1-p_i)\le\exp(-\sum_i p_i)\).
Finally \(\mu=0\) forces each \(p_i=0\), including the
possibility of an empty sum, proving the stated degeneracy
\cite{alon}.
\end{proof}
\chapter{The Lov\'asz Local Lemma}
\label{p205:local-lemma}

The finite union bound avoids all bad events when their total probability
is below \(1\). The local lemma allows a different test in which each
bad event is controlled by the number and probabilities of events
that can interact with it. We state the exact independence hypothesis
and prove the product lower bound for simultaneous avoidance.
Classical accounts appear in \cite[Section~5.1]{alon} and
\cite[Sections~2.2--2.3]{vondrak}.

\section{Dependency graphs}

\begin{definition}[A dependency graph for events]
\label{p205:W27}
Let \((A_i)_{i\in I}\) be a finite family of measurable events on a
probability space. A simple graph \(D\) on \(I\), with neighbor sets
\(N(i)\), is a dependency graph for the events if each \(A_i\) is
independent of the sigma-algebra generated by the events
\(A_j\) with \(j\notin N(i)\cup\{i\}\). Explicitly, for every event
\(B\) in that sigma-algebra one requires
\[
   \mathbb P(A_i\cap B)=\mathbb P(A_i)\mathbb P(B).
\]
\end{definition}

For example, suppose the experiment is a finite product of finite
weighted sample spaces, with independent coordinates. For each
\(i\) let \(A_i\) be determined by the coordinates in a set \(S_i\).
Join distinct \(i,j\) when \(S_i\cap S_j\ne\varnothing\). Each \(A_i\)
is a measurable cylinder event on its coordinate set. All the
non-neighbor events of \(A_i\) are measurable with respect to the
coordinates in \(T_i=\bigcup_{j\notin N(i)\cup\{i\}} S_j\).
Complements and countable unions of events determined by \(T_i\)
are still determined by \(T_i\), so the sigma-algebra generated by
the non-neighbor events has the same measurability property.
The sets \(S_i,T_i\) are disjoint. To see the required independence
directly, sum the product probabilities over assignments on the
two disjoint blocks and over the unused coordinates. The sum over
unused coordinates is \(1\), and the remaining sum factors as the
product of the two marginal probabilities. This also applies if a
block is empty.

Pairwise independence does not suffice for
Definition~\ref{p205:W27}. Let \(U,V\) be independent fair bits.
The events \(A=\{U=0\}\), \(B=\{V=0\}\), and \(C=\{U=V\}\)
each have probability \(1/2\) and each pair has intersection
probability \(1/4\). But \(A\cap B\) belongs to the sigma-algebra
generated by \(A,B\), and
\[
   \mathbb P(C\cap(A\cap B))=\tfrac14
      \ne\mathbb P(C)\mathbb P(A\cap B)=\tfrac18 .
\]
Thus the graph with no edges fails the dependency-graph condition
for this family.

\section{The asymmetric and symmetric bounds}

\begin{theorem}[Finite asymmetric Lov\'asz local lemma]
\label{p205:W28}
Let \(D\) be a dependency graph for a finite event family
\((A_i)_{i\in I}\). Suppose numbers \(x_i\in[0,1)\) satisfy
\[
   \mathbb P(A_i)\le x_i\prod_{j\in N(i)}(1-x_j)
   \qquad(i\in I).
\]
Then
\[
   \mathbb P\!\left(\bigcap_{i\in I}A_i^c\right)
     \ge \prod_{i\in I}(1-x_i)>0 .
\]
\end{theorem}

\begin{proof}
For \(S\subseteq I\), write \(B_S=\bigcap_{j\in S}A_j^c\) and
\(B_\varnothing=\Omega\). For a positive-probability event \(B\),
write \(\mathbb P(A\mid B)=\mathbb P(A\cap B)/\mathbb P(B)\).
We prove by induction on \(|S|\) both
\(\mathbb P(B_S)>0\) and
\begin{equation}\label{p205:lll-conditional}
     \mathbb P(A_i\mid B_S)\le x_i
     \qquad(i\notin S).
\end{equation}
For \(S=\varnothing\), positivity is immediate and the hypothesis
gives \(\mathbb P(A_i)\le x_i\). Suppose \(|S|=m\ge1\) and both
statements are established for all smaller sets. Select \(j\in S\).
By the conditional estimate at \(S\setminus\{j\}\),
\[
 \mathbb P(B_S)=\mathbb P(B_{S\setminus\{j\}})
         (1-\mathbb P(A_j\mid B_{S\setminus\{j\}}))
     \ge \mathbb P(B_{S\setminus\{j\}})(1-x_j)>0 .
\]

To establish \eqref{p205:lll-conditional} at this size, fix \(i\notin S\)
and write \(U=S\cap N(i)\) and \(T=S\setminus U\). The event \(B_T\)
is in the non-neighbor sigma-algebra of \(A_i\) and has positive
probability. Therefore
\[
  \mathbb P(A_i\mid B_S)=
    \frac{\mathbb P(A_i\cap B_U\mid B_T)}
         {\mathbb P(B_U\mid B_T)}
    \le \frac{\mathbb P(A_i)}{\mathbb P(B_U\mid B_T)}.
\]
If \(U\) is empty the denominator is \(1\). Otherwise enumerate
\(U\) as \(u_1,\ldots,u_h\). The conditional chain rule gives
\[
 \mathbb P(B_U\mid B_T)
 =\prod_{\ell=1}^h\!
     \left(1-\mathbb P\bigl(A_{u_\ell}\mid
            B_{T\cup\{u_1,\ldots,u_{\ell-1}\}}\bigr)\right)
 \ge \prod_{j\in U}(1-x_j).
\]
Every conditioning set in this product has size at most \(m-1\)
and excludes the event being estimated. In the product before the
inequality, each factor is the ratio of the probability after
adjoining \(u_\ell\) to the probability before it. The ratios
telescope to \(\mathbb P(B_S)/\mathbb P(B_T)\). All terms are
well-defined and the bounds follow from the induction hypothesis.
The same lower bound holds for the empty product when \(U\) is empty.
Using the hypothesis on \(\mathbb P(A_i)\) now gives
\[
  \mathbb P(A_i\mid B_S)\le
   x_i\prod_{j\in N(i)\setminus U}(1-x_j)\le x_i,
\]
completing the induction. Finally enumerate \(I\) and apply the
conditional chain rule to the complements in this order. Estimate
each factor by \eqref{p205:lll-conditional} to obtain the required
product. Every factor is strictly positive. In particular the argument
allows \(x_i=0\). The empty family gives the probability and product
both equal to \(1\).
\end{proof}

\begin{corollary}[Finite symmetric Lov\'asz local lemma]
\label{p205:W29}
Suppose a dependency graph for the event family has maximum degree
at most an integer \(d\ge0\), and that \(\mathbb P(A_i)\le p\) for
some \(p\ge0\) and every \(i\). If \(e\,p(d+1)\le1\), then
\(\mathbb P(\bigcap_i A_i^c)>0\).
\end{corollary}

\begin{proof}
If \(d\ge1\), use the same \(x=1/(d+1)\in(0,1)\) for every event.
Since \(\log(1+s)\le s\) for \(s\ge0\) by integrating
\(1/(1+s)\le1\), one has \((1+1/d)^d\le e\). Consequently
\[
 x(1-x)^{|N(i)|}\ge \frac{1}{d+1}\left(\frac d{d+1}\right)^d
     \ge\frac{1}{e(d+1)}\ge p .
\]
If \(d=0\), take \(x=p\) with \(0\le p\le1/e<1\), since all neighbour
products are \(1\). In either case apply Theorem~\ref{p205:W28}.
The empty event family has avoidance probability \(1\).
\end{proof}

\section{An application to hypergraph coloring}

\begin{theorem}[Property B under bounded edge intersections]
\label{p205:W30}
Let a finite hypergraph have every edge of size at least \(k\ge1\).
Suppose each edge intersects at most \(d\ge0\) other edges.
If \(e\,2^{1-k}(d+1)\le1\), there is a red--blue coloring
of the vertices with no monochromatic edge.
\end{theorem}

\begin{proof}
Assign each vertex an independent fair red--blue color.
For an edge \(F\), the two monochromatic patterns on \(F\) are
distinct because \(F\) is nonempty. The probability of the bad
event is therefore \(2^{1-|F|}\le 2^{1-k}\). Its coordinate set
is the vertex set of \(F\), so the dependency graph obtained by
joining intersecting edges satisfies Definition~\ref{p205:W27}
by the product argument above. Apply
Corollary~\ref{p205:W29} with \(p=2^{1-k}\) and take a witnessing
coloring by Proposition~\ref{p205:W14}. If the hypergraph has
no edges every coloring is a witness.
\end{proof}
\chapter{Linear Programming and Total Unimodularity}

\section{Duality and certificates}

Vectors and matrices in this chapter are finite and real unless
integrality is specified. Inequalities on vectors are coordinatewise.
We use the finite-dimensional linear programming theorem and Farkas
alternative proved in \cite[Theorems 5.1 and 1.23]{p106}. Schrijver
\cite[vol.~A, Sections 5.4--5.6 and 13.2]{schrijver} gives a parallel
treatment.

\begin{theorem}[Linear programming certificates]\label{p205:W31}
Consider the primal and dual
\[
 \inf\{c^{\mathsf T}x:Ax=b,\ x\ge0\},
 \qquad
 \sup\{b^{\mathsf T}y:A^{\mathsf T}y\le c\}.
\]
For feasible \(x,y\), their objective values are equal exactly when
\(x_i(c-A^{\mathsf T}y)_i=0\) for every \(i\). Either condition certifies
optimality on both sides. If the primal is feasible with finite
infimum, there is such a pair. The primal is infeasible exactly when
there exists \(u\) with \(A^{\mathsf T}u\ge0\) and \(b^{\mathsf T}u<0\).
For a feasible primal, the infimum is \(-\infty\) exactly when there
exists \(d\ge0\) with \(Ad=0\) and \(c^{\mathsf T}d<0\).
\end{theorem}

\begin{proof}
Subtract the objectives and use \(Ax=b\) to get
\(x^{\mathsf T}(c-A^{\mathsf T}y)\). Its summands are nonnegative.
This proves weak duality, the equality test and the optimality
certificate. Attainment of both optima and equality of their values
in the finite case are precisely \cite[Theorem 5.1]{p106} in this
standard form.

If both \(x\) and \(u\) in the infeasibility test existed, then
\(0\le x^{\mathsf T}A^{\mathsf T}u=b^{\mathsf T}u<0\). The converse
is the Farkas alternative for \(Ax=b,\ x\ge0\).
For the unboundedness test, start at any feasible \(x_0\) and use
\(x_0+td\) for \(t\ge0\). Conversely, unboundedness makes the dual
infeasible by weak duality. Write its feasibility conditions as
\[
 [\,A^{\mathsf T},-A^{\mathsf T},I\,](y^+,y^-,s)^{\mathsf T}=c,
 \qquad y^+,y^-,s\ge0 .
\]
The Farkas certificate for this system is a vector \(d\) with
\(Ad\ge0\), \(-Ad\ge0\), \(d\ge0\), and \(c^{\mathsf T}d<0\),
as required.
\end{proof}

We can pass to the standard form from a system with free variables,
mixed equalities and inequalities, and either objective direction.
Split each free coordinate into the difference of two nonnegative
ones, add one slack to each upper inequality, and negate the
objective for maximization. These transformations preserve the set
of objective values. A row used as an upper inequality in a
maximization has a nonnegative dual multiplier because its slack
has objective coefficient zero. An equality has a free multiplier.
We apply this correspondence explicitly to the two matching models
in Chapter~\ref{p205:ch08}.

\section{Vertices and active rows}

A point \(x\) of a convex set is extreme if it admits no expression
\(x=(x_++x_-)/2\) with two distinct points \(x_+,x_-\) of the set.
This midpoint test also rules out any nontrivial convex
decomposition into two points, by taking a smaller symmetric
segment within the latter segment. A polyhedron is pointed if it
contains no affine line. We use vertex to mean extreme point of a
polyhedron.

\begin{theorem}[Active rows and vertex optima]\label{p205:W32}
Let \(P=\{x\in\mathbb R^n:Ax\le b\}\). A point \(x\in P\) is a vertex
exactly when the rows with \(a_i^{\mathsf T}x=b_i\) span the dual
coordinate space. If \(P\) is nonempty and pointed, a linear
functional with finite optimum on \(P\) attains it at a vertex.
Every nonempty bounded \(P\) is the convex hull of its vertices.
\end{theorem}

\begin{proof}
If the active rows fail to span, choose a nonzero \(d\) annihilated
by all of them. The finitely many inactive inequalities have
strict slack, so \(x\pm\epsilon d\) are feasible for sufficiently
small \(\epsilon>0\). Conversely, if \(x\) is the midpoint of two
feasible points, every row active at \(x\) is active at both
endpoints. Spanning forces the endpoints to coincide.

Start at an optimal point, whose existence follows from the
standard-form transformation and Theorem~\ref{p205:W31}.
If the active rows do not span, take \(d\) as above. A small step
in either direction is feasible, so the objective vanishes on \(d\).
If every row vanished on \(d\), the entire line through \(x\)
parallel to \(d\) would be feasible. By pointedness, change the
sign of \(d\) if necessary to make some row positive on \(d\).
The first new active constraint occurs after a step of length
\[
 t_*=\min_{i:a_i^{\mathsf T}d>0}
          \frac{b_i-a_i^{\mathsf T}x}{a_i^{\mathsf T}d}.
\]
The minimum is over a nonempty finite set and is positive. The new
row is independent of the old active span, since it is positive
on \(d\). The rank therefore increases without changing the
objective. After finitely many steps the active rows span.

For the last assertion start at any feasible \(x\). If the active
rows fail to span, use a nonzero \(d\) in their common kernel.
Boundedness gives first hits in both directions, at positive
distances \(t_+,t_-\). The point \(x\) is the convex combination of
\(x+t_+d\) and \(x-t_-d\) with respective weights
\(t_-/(t_++t_-)\) and \(t_+/(t_++t_-)\). At both new points
the active rank increases. Induction on the rank deficit expresses
\(x\) as a finite convex combination of vertices. The reverse
containment is convexity.
\end{proof}

\section{Integer hulls and incidence matrices}

\begin{definition}[Total unimodularity]\label{p205:W33}
An integer matrix is \emph{totally unimodular} if every square
submatrix has determinant in \(\{0,1,-1\}\). We include the
empty determinant \(1\).
\end{definition}

\begin{theorem}[Integer right-hand sides]\label{p205:integer-hull}
Let \(A\) be totally unimodular and \(b\) integer.
Every vertex of \(P=\{x:Ax\le b\}\) is integral. In fact
\(P=\operatorname{conv}(P\cap\mathbb Z^n)\), with the empty
hull interpreted as empty. For nonempty \(P\), a finite
linear optimum has an integer optimizer if \(P\) is pointed,
in particular if it is bounded.
\end{theorem}

\begin{proof}
First note that transposition, submatrices, row sign changes,
and row repetition preserve total unimodularity. For repetition,
a square minor with two identical rows vanishes. Appending rows
from \(I\) and \(-I\) also preserves the property, since expansion
along any such row reduces a nonzero minor to a smaller one up to
sign. These facts cover bounds, and cover replacing an equality
row by two opposite inequalities.

At a vertex take \(n\) independent active rows by
Theorem~\ref{p205:W32}. Their square matrix has determinant
\(\pm1\), so Cramer's rule expresses every coordinate as an integer.
Now for any \(x\in P\), add the coordinate box
\(\lfloor x_i\rfloor\le z_i\le\lceil x_i\rceil\) on the variable
\(z\). This is a nonempty bounded polyhedron with integer right-hand
side and a totally unimodular constraint matrix. Its vertices are
integral and lie in \(P\), and its vertex hull contains \(x\)
by Theorem~\ref{p205:W32}. This argument also works if \(P\) contains
a line. The reverse hull inclusion follows from convexity. Finally
a finite optimum on pointed \(P\) has an optimal vertex by
Theorem~\ref{p205:W32}.
\end{proof}

\begin{theorem}[Hoffman--Kruskal criterion]\label{p205:W34}
An integer matrix \(A\) is totally unimodular exactly when every
polyhedron \(\{x:x\ge0,\ Ax\le b\}\) with integral \(b\) has
integral vertices.
\end{theorem}

\begin{proof}
The forward direction follows by adjoining \(-I\) and applying
Theorem~\ref{p205:integer-hull}. For the converse, suppose that
a nonsingular submatrix \(B\) on row set \(I_0\) and column set
\(J_0\) has determinant of absolute value at least two. There is
an integer vector \(q\) with \(B^{-1}q\) nonintegral. Indeed, if
every column \(B^{-1}e_i\) were integral, the determinant of
\(B^{-1}\) would be an integer. Choose an integer vector \(k\)
so that \(z=B^{-1}q+k\) has positive entries, and set
\(x_{J_0}=z\), \(x_{J_0^c}=0\). On \(I_0\) use the integer right-hand
side \(q+Bk\), and on all other rows use the ceilings of the
row products \(a_i^{\mathsf T}x\). The resulting point is feasible,
with \(I_0\) and all the coordinate rows on \(J_0^c\) active.
A vector annihilated by these rows must first vanish on \(J_0^c\)
and then vanish on \(J_0\) by invertibility of \(B\). Thus \(x\)
is a vertex, contradicting integrality.
\end{proof}

\begin{theorem}[Network and bipartite integrality]\label{p205:W35}
The node--arc incidence matrix of a finite directed network and the
unsigned vertex--edge incidence matrix of a finite bipartite graph
are totally unimodular. If \(B\) is a node--arc incidence matrix
using the net-outflow sign, \(b,\ell,u\) are integral and the bounds
are finite, all vertices of
\[
             \{f:Bf=b,\ \ell\le f\le u\}
\]
are integral. A nonempty such polyhedron has an integral optimizer
for every linear objective. For a bipartite graph \((V,E)\), the
matching polytope is exactly
\[
 \left\{x\in\mathbb R^E:x\ge0,\
              \sum_{e\ni v}x_e\le1\ \text{for every }v\in V\right\}.
\]
Replacing every degree inequality by equality gives the hull of
perfect matchings. Here a matching polytope denotes the convex
hull of the respective incidence vectors.
\end{theorem}

\begin{proof}
A nonloop arc contributes \(+1\) at its tail and \(-1\) at its
head. A loop has zero column. In any square submatrix, a column
with at most one nonzero entry allows determinant expansion.
If every column has two nonzero entries, the row sum is zero,
and the determinant vanishes. Induction gives total unimodularity,
also for parallel arcs. For a bipartite graph, negating the rows
of one part gives the directed incidence matrix of the
orientation from the other part.

Equality and bound constraints preserve the minor property by
the closure rules in Theorem~\ref{p205:integer-hull}. This gives
the flow assertions. In the bipartite systems each nonnegative
edge coordinate is at most one by either of its incident vertex
constraints. The polyhedra are bounded, so each is the hull of
its integral vertices. Integral feasible points are precisely
the required zero-one matching vectors, or perfect matching
vectors in the equality case. If there is no perfect matching,
the latter polyhedron is empty, because otherwise it has an
integral vertex.
\end{proof}

\begin{example}[The degree relaxation on a triangle]
Set each edge coordinate to \(1/2\). The three degree inequalities
are equalities. On a fixed ordering of the edges the active
equations are \(x_1+x_2=x_2+x_3=x_3+x_1=1\), with the unique
solution \((1/2,1/2,1/2)\). Thus this is a fractional vertex.
Every actual matching in the triangle has size at most one.
\end{example}
\chapter{Network Flows and Assignment}\label{p205:ch08}

\section{Residual networks and shortest augmentation}

A finite directed network here has a finite vertex set \(V\) and
a set of labelled arcs \(E\), with a tail and a head for each arc.
Parallel and oppositely oriented arcs can have different labels.
A flow is a real vector indexed by arcs. Its net outflow at \(v\)
is the sum on outgoing arcs minus the sum on incoming arcs.
A loop contributes zero to net outflow. Write \(\delta^+(S)\) and
\(\delta^-(S)\) for the arcs leaving and entering \(S\), respectively.
For an arc vector \(a\) and an arc set \(F\), put
\(a(F)=\sum_{e\in F}a_e\).

For bounds \(\ell\le f\le u\), the forward residual copy of an
arc has capacity \(u_e-f_e\) and the backward copy has capacity
\(f_e-\ell_e\), each present when its capacity is positive. These
copies retain the original arc label and a direction sign.
Augmentation adds the selected amount along forward copies and
subtracts it along backward copies. A simple residual path is
limited by its minimum residual capacity. A simple residual
cycle can be augmented by the same rule. If both copies of one
arc occur on a cycle their increments cancel, so feasibility
still follows from the individual residual bounds.

\begin{theorem}[Flow--cut optimality certificates]\label{p205:W36}
Let \(s,t\) be distinct vertices and \(u_e\ge0\) be finite real
capacities. Consider flows \(0\le f\le u\) conserving at vertices
other than \(s,t\), with value equal to the net outflow at \(s\).
A flow and a cut \(S\) with \(s\in S,\ t\notin S\) certify each
other's optimality when the flow value equals \(u(\delta^+(S))\).
There always exists a pair of equal value. If all capacities
are integers, the flow in such a pair can be taken integral.
\end{theorem}

\begin{proof}
Summing the net outflows over \(S\) gives
\[
     \operatorname{val}(f)=f(\delta^+(S))-f(\delta^-(S))
                            \le u(\delta^+(S)).
\]
This proves the certificate. With real capacities, the flow
polyhedron is closed, bounded and nonempty because it contains
zero. Take a maximizing flow. A residual \(s\)--\(t\) path would
improve it, so the set reachable from \(s\) is a cut. Every
outgoing arc is saturated and every incoming arc carries zero,
since otherwise the reachability set would grow. Thus its
capacity is the flow value. This is the same real-capacity
argument as in \cite[Chapter 7, W44--W46]{p204}.

For integral \(u\), choose an integral optimal vertex by
Theorem~\ref{p205:W35}. To apply that theorem, adjoin an arc
\(t\to s\) with bounds \([-M,M]\), where
\(M=\sum_{e\in E}u_e\). Set its flow to the value of the
original flow and impose zero net outflow at every vertex.
This gives a bijective extension of the original feasible set
because the value has absolute value at most \(M\). Maximizing
the new arc coordinate gives an integral original maximizer.
The reachable cut construction still applies.
\end{proof}

\begin{theorem}[Shortest augmenting paths]\label{p205:W37}
Starting from zero in the setting of Theorem~\ref{p205:W36},
augment by the full residual bottleneck along a shortest
\(s\)--\(t\) residual path as long as there is one. With real
capacities, this returns a maximum flow after \(O(|V||E|)\)
augmentations. A breadth-first search and update per iteration
use \(O(|V|+|E|)\) arc inspections and elementary operations.
\end{theorem}

\begin{proof}
Loops can be ignored. Let \(d(v)\) and \(d'(v)\) be the distances
from \(s\) before and after an augmentation, with distance
\(\infty\) for an unreachable vertex. On the selected shortest
path every edge goes from distance \(j\) to \(j+1\), because its
prefixes are also shortest. Suppose some distance decreases, and
follow a new shortest path to the first vertex \(v\) with
\(d'(v)<d(v)\), reached by an arc \(u\to v\). Its predecessor
satisfies \(d(u)\le d'(u)<\infty\). If the arc was present before
the augmentation then \(d(v)\le d(u)+1\le d'(v)\), a contradiction.
A newly available arc is the reverse copy of an arc \(v\to u\)
on the selected path. In this case \(d(v)=d(u)-1<d'(v)\), again
a contradiction. Hence distances never decrease.

In each augmentation some residual copy is saturated and
disappears. If the same labelled copy \(u\to v\) disappears
by saturation again, its paired copy \(v\to u\) must have been
used in between to restore its capacity. At the earlier
saturation \(d(v)=d(u)+1\). At the intervening use of \(v\to u\),
the then-current distance of \(u\) is the distance of \(v\) plus
one, so it is at least the earlier distance of \(u\) plus two.
Distances of reachable vertices on a simple path are bounded by
\(|V|-1\). Thus any fixed labelled copy can be charged at most
\(O(|V|)\) times. There are at most \(2|E|\) copies, so only
\(O(|V||E|)\) augmentations are possible. Searching adjacency
lists, tracing the path and updating its copies give the stated
per-iteration count, with one final search if no path remains.
The cut of reachable vertices then certifies optimality.
\end{proof}

This is the Edmonds--Karp shortest-path rule in
\cite[vol.~A, Section 10.5]{schrijver}. The operation bound
counts exact residual operations, independently of the values
of the capacities.

\section{Bounded flows and minimum costs}

\begin{theorem}[Hoffman circulation and fixed divergence]\label{p205:W38}
Let \(\ell_e\le u_e\) be finite real arc bounds and let \(b\)
be a real vertex vector with \(\sum_v b_v=0\). There is a flow
within the bounds with net outflow \(b\) exactly when
\[
 b(S)\le u(\delta^+(S))-\ell(\delta^-(S))
                       \qquad(S\subseteq V).
\]
Here \(b(S)=\sum_{v\in S}b_v\). If \(b,\ell,u\) are integral,
there is an integral feasible flow whenever these conditions hold.
The case \(b=0\) gives the circulation criterion.
\end{theorem}

\begin{proof}
Summing the required net outflows over \(S\) gives necessity.
For sufficiency, let \(B\) denote the net-outflow incidence
matrix and write \(g=f-\ell\), \(c=u-\ell\),
\(\beta=b-B\ell\). We need \(0\le g\le c\) and \(Bg=\beta\).
The cut inequalities become
\(\beta(S)\le c(\delta^+(S))\), and \(\sum_v\beta_v=0\).

Add a new source \(s_*\), a new sink \(t_*\), and arcs
\(s_*\to v\) of capacity \(\beta_v\) for \(\beta_v>0\) and
\(v\to t_*\) of capacity \(-\beta_v\) for \(\beta_v<0\).
Let \(D=\sum_{\beta_v>0}\beta_v=\sum_{\beta_v<0}(-\beta_v)\).
A source--sink flow of value \(D\) saturates every artificial
arc, so its restriction has net outflow \(\beta\). Conversely
a solution \(g\) extends to a flow of value \(D\) using precisely
those artificial loads. For a cut with original part \(S\), the
capacity is
\[
 c(\delta^+(S))
     +\sum_{v\notin S,\ \beta_v>0}\beta_v
     +\sum_{v\in S,\ \beta_v<0}(-\beta_v)
       =D+c(\delta^+(S))-\beta(S).
\]
Thus all cuts have capacity at least \(D\).
Theorem~\ref{p205:W36} produces the required flow, integrally
when the data are integral. Add \(\ell\) back. Empty sets and
\(D=0\) fit the same formulas.
\end{proof}

\begin{theorem}[Minimum-cost flow and potential certificates]\label{p205:W39}
Let \(f\) be feasible for the bounded fixed-divergence problem of
Theorem~\ref{p205:W38}. Give each arc a real cost \(w_e\), with
residual costs \(w_e\) forward and \(-w_e\) backward. Then the
following are equivalent.
\begin{enumerate}
\item The flow \(f\) minimizes \(\sum_e w_ef_e\).
\item Its residual graph has no directed cycle of strictly negative
total cost, with loops allowed as cycles.
\item There is a real vertex function \(\pi\) such that every
residual arc \(v\to z\) of cost \(w\) satisfies
\(w+\pi(v)-\pi(z)\ge0\).
\end{enumerate}
For integral bounds and divergence, a nonempty feasible set admits
an integral minimum-cost flow, also for real costs.
\end{theorem}

\begin{proof}
A negative-cost cycle gives a feasible improving augmentation by
its positive bottleneck. Both copies of the same original arc
would form a zero-cost two-arc cycle if they comprised the entire
cycle, and in any event their signed contributions cancel.
This proves the first implication. For the reverse implication,
compare \(f\) with another feasible \(h\). Put the positive parts
of \(h_e-f_e\) on forward copies and the negative parts with
opposite sign on backward copies. The resulting nonnegative
residual arc vector is a circulation, since \(h\) and \(f\)
have the same divergence. Decompose it into directed cycles
with positive coefficients as in \cite[Chapter 7, W47]{p204}.
Concretely, follow positive arcs until a vertex repeats, subtract
the minimum load on the resulting cycle, and repeat. Conservation
persists and the positive support shrinks each time. The cost
difference is the same weighted sum of cycle costs, so it is
nonnegative if there is no negative cycle.

To construct \(\pi\), add a new vertex with a zero-cost outgoing
arc to each original vertex and no incoming arc. Let \(\pi(v)\)
be the minimum cost of a path from the new vertex to \(v\).
Every walk can have its closed subwalks removed without increasing
its cost. Indeed, a closed walk decomposes into directed cycles
by splitting at repeated vertices, and each has nonnegative cost.
Thus the minimum over the finitely many simple paths exists
and equals the infimum over walks. Appending an arc of cost
\(w\) gives \(\pi(z)\le\pi(v)+w\), as claimed. Conversely,
sum the reduced-cost inequalities around a cycle and cancel
the potentials. The final integrality assertion is
Theorem~\ref{p205:W35} applied to the cost objective.
\end{proof}

The feasible-flow and reduced-cost constructions follow the
interfaces in \cite[vol.~A, Sections 11.2--11.4 and 12.1--12.2]{schrijver}.

\section{Bipartite assignment and its dual}

Let \(L\) and \(R\) be the two sides of a finite bipartite graph.
An allowed edge \(ij\) joins \(i\in L\) to \(j\in R\), with
real weight \(w_{ij}\). The incidence matrix on the allowed edges
is denoted by \(A\), with rows indexed by the vertices. The
following two forms use different sign conditions for their
dual variables.

\begin{theorem}[Perfect assignment duality]\label{p205:W40}
Assume that a perfect matching exists. Its maximum possible weight
equals
\[
   \min\left\{\sum_{i\in L}u_i+\sum_{j\in R}v_j:
          u_i+v_j\ge w_{ij}\ \text{for each allowed }ij\right\},
\]
where the potentials are free real numbers. The minimum is attained.
For feasible potentials a perfect matching is jointly optimal with
them exactly when all its edges are tight.
\end{theorem}

\begin{proof}
The linear relaxation maximizes \(w^{\mathsf T}x\) under \(Ax=\mathbf
1,\ x\ge0\). By Theorem~\ref{p205:W35} it has a perfect matching
as an optimizer. Minimize \(-w^{\mathsf T}x\) instead and apply
Theorem~\ref{p205:W31}. Negating the equality multiplier yields
the displayed minimization with no sign restriction. The two
objectives agree at a joint optimum. In general the difference
between potential sum and matching weight is the sum, over the
matching, of nonnegative edge slacks, giving the tightness test.
\end{proof}

\begin{theorem}[Weighted bipartite matching duality]\label{p205:W41}
Without a perfect matching requirement, the maximum matching weight
equals the same minimum with \(u_i,v_j\ge0\). The empty matching is
allowed and the minimum is attained. A matching and feasible
potentials certify each other's optimality precisely when the
matching uses tight edges and all its unmatched vertices have
zero potential.
\end{theorem}

\begin{proof}
The relaxation is now \(Ax+s=\mathbf 1,\ x,s\ge0\), with objective
\(w^{\mathsf T}x\). Theorem~\ref{p205:W35} identifies an optimizer
that is a matching. Apply Theorem~\ref{p205:W31} to the minimization
of its negative. The constraints on the slack columns give
nonnegativity of the negated dual multipliers, and the edge columns
give \(u_i+v_j\ge w_{ij}\). The difference of the objectives is
the sum of edge slacks on the matching plus the potentials of the
unmatched vertices. All these terms are nonnegative, which proves
the certificate in both directions.
\end{proof}

\begin{theorem}[Birkhoff--von Neumann decomposition]\label{p205:W42}
A finite square matrix with nonnegative entries and every row and
column sum equal to one is a convex combination of permutation
matrices. Its extreme points, in the set of all such matrices of the
same size, are precisely the permutation matrices. If some positions
are prescribed to be zero, all permutations in the decomposition
can be chosen supported on the allowed positions. These permutation
matrices are also precisely the extreme points of the restricted set.
\end{theorem}

\begin{proof}
The zero-size matrix is itself the empty permutation matrix.
For positive size, build the bipartite graph of positive entries.
For any set \(S\) of rows, its total mass \(\#S\) lies in its
neighbor columns \(N(S)\), with each receiving at most one.
Thus \(\#S\le\#N(S)\), and Hall's theorem
\cite[Chapter 8, W51]{p204} gives a perfect matching.
Let \(P\) be its permutation matrix, and let \(\lambda>0\)
be the smallest matrix entry on that matching. If \(\lambda=1\),
all other entries vanish by the row sums. Otherwise \(X-\lambda P\)
is nonnegative, has row and column sums \(1-\lambda\), and has
strictly smaller positive support. Induct on support after dividing
by \(1-\lambda\) and use
\[
             X=\lambda P+(1-\lambda)\frac{X-\lambda P}{1-\lambda}.
\]
All permutations respect the original positive support. Conversely
every permutation matrix is extreme, since in a convex
decomposition its zero positions force zeros in both summands
and row sums then force the ones. A convex combination involving
distinct permutations is not extreme, as follows by isolating
one of its positive coefficients. The same argument applies with
prescribed zero positions. The unrestricted theorem is also
available in the fixed Mathlib \cite{mathlib}.
\end{proof}

\begin{theorem}[Hungarian method on allowed edges]\label{p205:W43}
For a finite bipartite graph that has a perfect matching, the
tight-edge Hungarian method terminates after finitely many
potential updates and matching augmentations. It returns a
maximum-weight perfect matching and an optimal feasible potential
pair for Theorem~\ref{p205:W40}.
\end{theorem}

\begin{proof}
The empty graph with empty sides is immediate. Otherwise both sides
have the same cardinality, and each left vertex has an allowed
neighbor. Start with \(v_j=0\) and
\(u_i=\max_{j:ij\text{ allowed}}w_{ij}\), and an empty matching.
Maintain feasible potentials and a matching on tight edges.
Fix an unmatched left vertex \(i_0\). Search from \(i_0\) along
tight edges outside the matching from left to right, and along
matched edges from right to left. If this reaches an unmatched
right vertex, a simple alternating path gives a tight-edge
augmentation.

Otherwise let \(S\) and \(T\) be the reached left and right vertices.
Each reached right vertex is matched to a reached left vertex, and
each reached left vertex other than \(i_0\) arrived by such a
matched edge. Hence \(\#S=\#T+1\). A perfect matching in the allowed
graph requires an allowed edge from \(S\) to \(R\setminus T\),
since its edges from \(S\) have distinct right endpoints.
No such edge is matched in the current matching and no such edge
is tight, by the search. Take the minimum \(\delta>0\) of
\(u_i+v_j-w_{ij}\) over these edges. Subtract \(\delta\) from
\(u_i\) for \(i\in S\) and add it to \(v_j\) for \(j\in T\).
Slacks on \(S\times(R\setminus T)\) decrease by at most their
old value. Those on \((L\setminus S)\times T\) increase and all
others stay fixed. Matched edges have both endpoints inside the
reached sets or both outside, so they stay tight.
The previously reached alternating paths also stay tight.
At least one new right vertex becomes reachable via a newly tight
edge from \(S\). Thus for this root, repeated updates either reach an
unmatched right vertex or strictly enlarge the reached right set.
They must eventually give an augmentation, since an unsuccessful
search always has \(\#T<\#L=\#R\) and an edge available for the next
update. Each augmentation increases matching size by one. The
resulting perfect matching is tight and certifies optimality by
Theorem~\ref{p205:W40}.
\end{proof}

See Kuhn \cite{kuhn} and \cite[vol.~A, Sections 17.2--17.3]{schrijver}
for the assignment method and its historical form.
\chapter{Matroid Polytopes and Optimization}
\label{p205:ch09}

A finite matroid on \(E\) is a family of independent subsets
containing the empty set, closed under subsets, and satisfying
augmentation. This last axiom says that for independent \(I,J\)
with \(|I|<|J|\), some \(e\in J\setminus I\) has \(I\cup\{e\}\)
independent. Write \(r(S)\) for the maximum size of an independent
subset of \(S\). A basis of \(S\) is a maximal independent subset of
\(S\). By augmentation all such bases have size \(r(S)\). The bases
of the matroid are the bases of \(E\).
This is the convention of \cite[Chapter~10, W64]{p204}.

For \(x\in\mathbb R^E\) and \(S\subseteq E\), write
\(x(S)=\sum_{e\in S}x_e\), and write \(\mathbf 1_S\) for the
indicator vector of \(S\). The dot product is
\(w\cdot x=\sum_{e\in E}w_e x_e\).

\section{Greedy weights and rank polytopes}

\begin{theorem}[Maximum-weight basis]
\label{p205:W44}
Assign arbitrary real weights to a finite matroid's elements.
Scan them in nonincreasing weight order, adding each element
whose addition preserves independence. Every tie order gives
a maximum-weight basis.
\end{theorem}

\begin{proof}
The empty ground set has the unique basis \(\varnothing\).
Otherwise list the elements as \(e_1,\ldots,e_n\), with
\(w_i=w_{e_i}\) and \(P_i=\{e_1,\ldots,e_i\}\).
Write \(B\) for the greedy output.
If an element of \(P_i\) was rejected, then it was dependent
when added to the set selected at its step. By heredity its
addition to \(B\cap P_i\) is still dependent. Thus \(B\cap P_i\)
is maximal independent in \(P_i\) and has size \(r(P_i)\).
In particular \(B\) is a basis. For every \(x\in\mathbb R^E\)
the telescoping identity is
\begin{equation}
\label{p205:matroid-prefix-sum}
w\cdot x=w_n x(E)+\sum_{i=1}^{n-1}(w_i-w_{i+1})x(P_i).
\end{equation}
Indeed the coefficient of \(x_{e_j}\) on the right is
\(w_n+\sum_{i=j}^{n-1}(w_i-w_{i+1})=w_j\).
For every basis \(D\) we have \(|D|=r(E)\) and
\(|D\cap P_i|\le r(P_i)\). The term involving \(w_n\) is constant
on bases even when \(w_n<0\), and every other coefficient is
nonnegative. The greedy basis attains all prefix bounds.
The same argument applies to any ordering within a tie block.
\end{proof}

For the wider greedy characterization of matroids see
\cite[Chapter~12, W78--W80]{p204}. The weight and polytope
arguments follow \cite[Sections~40.1--40.2]{schrijver} and
\cite[Section~1.8]{oxley}.

\begin{lemma}[A finite convex hull is determined by its linear bounds]
\label{p205:matroid-support}
Let \(F\) be a nonempty finite subset of a Euclidean space. If
\(w\cdot x\le\max_{z\in F} w\cdot z\) for every \(w\), then
\(x\in\operatorname{conv}F\).
\end{lemma}

\begin{proof}
Put \(K=\operatorname{conv}F\). This set is compact because it is
the image of the compact simplex of convex coefficients under a
linear map. If \(x\notin K\), minimize distance from \(x\) at some
\(y\in K\). For \(z\in K\) the squared distance from \(x\) to
\(y+t(z-y)\), \(0\le t\le 1\), is
\[
 \|x-y\|^2-2t(x-y)\cdot(z-y)+t^2\|z-y\|^2.
\]
Minimality at \(t=0\) gives \((x-y)\cdot(z-y)\le0\) by dividing
by positive \(t\) and letting \(t\) decrease to zero. With
\(w=x-y\) we have \(w\cdot z\le w\cdot y\) throughout \(K\), while
\(w\cdot x=w\cdot y+\|w\|^2>w\cdot y\). This contradicts the
assumed linear bound. For the general separation theorem see
\cite[Chapter~1, Open and compact--closed separation]{p106}.
\end{proof}

\begin{theorem}[Matroid base polytope]
\label{p205:W45}
For every finite matroid \(M\) with rank \(r\),
\[
 \begin{aligned}
 &\operatorname{conv}\{\mathbf 1_B\mid B\text{ is a basis of }M\}\\
 &\quad =
 \{x\in\mathbb R^E\mid x\ge 0,\ x(S)\le r(S)\text{ for all }S\subseteq E,\
                  x(E)=r(E)\}.
 \end{aligned}
\]
Its vertices are precisely the basis indicators and are integral.
\end{theorem}

\begin{proof}
Each basis indicator satisfies the constraints. The constraints
are linear, so their solution set contains the indicated hull.
For the reverse inclusion let \(x\) satisfy the constraints.
When \(E=\varnothing\) both sets consist of the unique empty vector.
Otherwise let \(w\in\mathbb R^E\) and order \(E\) as in
Theorem~\ref{p205:W44}. The greedy basis \(B\) has
\(\mathbf 1_B(P_i)=r(P_i)\) for each \(i\) by its prefix invariant.
In \eqref{p205:matroid-prefix-sum}, the term \(w_n x(E)\) equals
\(w_n\mathbf 1_B(E)\) and all other coefficients are nonnegative.
Hence \(w\cdot x\le w\cdot\mathbf 1_B\) for every \(w\).
Lemma~\ref{p205:matroid-support} gives the reverse inclusion.

An extreme point of a hull of finitely many points belongs to
the generating set, since an expression with positive mass on
distinct points would exhibit a proper convex decomposition.
Conversely, if a basis indicator is a convex combination of other
indicators, then each coordinate equal to \(0\) forces the same
coordinate to be \(0\) in every positive-weight summand.
A coordinate equal to \(1\) forces that coordinate to be \(1\)
in every positive-weight summand. Thus all such summands equal
the original indicator, proving extremality.
\end{proof}

\begin{theorem}[Matroid independence polytope]
\label{p205:W46}
For every finite matroid \(M\) with rank \(r\),
\[
 \operatorname{conv}\{\mathbf 1_I\mid I\text{ is independent in }M\}
 =
 \{x\in\mathbb R^E\mid x\ge 0,\ x(S)\le r(S)\text{ for all }S\subseteq E\}.
\]
For a finite undirected graph \(G=(V,E)\) the corresponding forest
polytope is
\begin{equation}
\label{p205:forest-polytope}
 \{x\in\mathbb R^E\mid
      x\ge0,\ x(E(S))\le |S|-1
      \text{ for all nonempty }S\subseteq V\},
\end{equation}
where \(E(S)\) is the set of edges with both endpoints in \(S\).
The spanning forest polytope adds \(x(E)=|V|-c(G)\), where
\(c(G)\) counts all components, including isolated vertices.
For a connected graph with at least one vertex this specializes
to the spanning tree polytope with \(x(E)=|V|-1\).
\end{theorem}

\begin{proof}
Every independent indicator satisfies the rank constraints.
For the reverse inclusion fix a solution \(x\) and a direction \(w\).
If \(w\) has no positive coordinate, \(w\cdot x\le0=w\cdot
\mathbf 1_\varnothing\). Otherwise list the positive elements
\(e_1,\ldots,e_m\) in nonincreasing order and put \(w_{m+1}=0\).
For \(P_i=\{e_1,\ldots,e_i\}\), telescoping and nonnegativity of \(x\)
give
\[
 w\cdot x\le \sum_{i=1}^m (w_i-w_{i+1})x(P_i)
 \le \sum_{i=1}^m (w_i-w_{i+1})r(P_i).
\]
The last quantity is \(w\cdot\mathbf 1_I\), where \(I\) is the
greedy selection on the positive elements. Its prefix ranks follow
from the proof of Theorem~\ref{p205:W44}. No nonpositive element is
needed to realize the bound. Now apply Lemma~\ref{p205:matroid-support}.

The graphic matroid has forests for independent sets, and its rank
on \(A\subseteq E\) is \(|V|-c(V,A)\), where \(c(V,A)\) counts
components of the spanning subgraph \((V,A)\).
Indeed a maximal forest in each component must span it. Otherwise
an edge joining two of the forest components could be added without
creating a cycle. A tree on \(q\ge1\) vertices has \(q-1\) edges,
as follows by repeatedly deleting a leaf. Summing over components
gives the rank formula. The matroid construction is the graphic
case of \cite[Chapter~11, W72]{p204}.

For nonempty \(S\), any forest of \(E(S)\) has at most \(|S|-1\)
edges, so the rank constraints imply \eqref{p205:forest-polytope}.
Conversely, for any \(A\subseteq E\) let \(S_1,\ldots,S_h\) be
the vertex sets of the components of \((V,A)\). Every edge of
\(A\) lies within one \(S_i\), and by \(x\ge0\) we get
\[
 x(A)\le\sum_{i=1}^h x(E(S_i))
      \le\sum_{i=1}^h(|S_i|-1)
      =|V|-h=r(A).
\]
The empty graph gives the empty sum. The spanning forests are
the graphic matroid bases, so the remaining statements follow
from Theorem~\ref{p205:W45}.
\end{proof}

\section{Spanning trees and exchange}

\begin{theorem}[Kruskal's algorithm]
\label{p205:W47}
In a finite undirected graph with real edge weights, process
edges in nondecreasing weight order and include an edge whenever
it does not create a cycle. For any ordering of tied edges, the
output is a spanning forest of minimum total weight. On a connected
graph with at least one vertex it is a minimum spanning tree.
\end{theorem}

\begin{proof}
A spanning forest consists of a spanning tree in every component
of the original graph. These are precisely the bases of the graphic
matroid. Processing edges by nondecreasing weight is basis greedy
with the negatives of those weights. Theorem~\ref{p205:W44}
therefore gives a minimum-weight basis. In the empty graph the
output is the empty forest.
\end{proof}

\begin{theorem}[Cut and cycle certificates and Prim's algorithm]
\label{p205:W48}
Let \(G\) be a finite connected undirected graph with at least one
vertex, with arbitrary real edge weights. Loops may be discarded,
and parallel edges are distinguished.
\begin{enumerate}
\item Let \(F\) be a forest contained in some minimum spanning tree.
If a nontrivial vertex cut has no edge of \(F\) crossing it, then
any minimum-weight edge crossing it can be added to \(F\) while
preserving extendibility to a minimum spanning tree.
Consequently Prim's algorithm gives a minimum spanning tree from
any starting vertex and with arbitrary ties.
\item A spanning tree \(T\) is minimum if and only if every
edge \(e\notin T\) with distinct endpoints satisfies
\[
 w_e\ge\max_{f\in P_T(e)}w_f,
\]
where \(P_T(e)\) is the unique \(T\)-path between those endpoints.
\item A spanning tree \(T\) is the unique minimum spanning tree if
and only if these inequalities are all strict.
\end{enumerate}
\end{theorem}

\begin{proof}
For the first assertion choose a minimum tree \(D\) containing \(F\)
and a light crossing edge \(e\). If \(e\notin D\), the unique path
in \(D\) between the endpoints of \(e\) crosses the cut, so choose
a crossing edge \(f\) on that path. By the cut condition \(f\notin F\).
Adding \(e\) and deleting \(f\) gives a tree containing \(F\cup\{e\}\).
Its weight is at most that of \(D\) because \(w_e\le w_f\).
It is therefore minimum. If \(e\in D\) the original \(D\) suffices.
Prim begins with a singleton visited set and no chosen edges.
When the visited set is proper, connectedness supplies a crossing
edge. Choosing a light one adds a new vertex and keeps the chosen
edges a tree on the visited set, so none of them crosses the next
cut. Induction applies the cut assertion at each step. At the
last step the chosen tree itself is minimum.

For the second assertion, if \(w_e<w_f\) for some \(f\in P_T(e)\),
adding \(e\) to \(T\) and removing \(f\) lowers the weight. Conversely
suppose the stated inequalities hold. For any spanning tree \(D\ne T\)
choose \(e\in D\setminus T\) and remove it. The endpoints of \(e\)
lie in the two different components of \(D-e\). The path \(P_T(e)\)
therefore has an edge \(f\) crossing the cut between the components.
The only edge of \(D\) crossing that cut was \(e\), so \(f\notin D\).
Replacing \(e\) by \(f\) gives another tree and changes its weight
by \(w_f-w_e\le0\). The replacement decreases \(|D\setminus T|\)
by one. Iteration ends at \(T\), whose weight is thus at most
that of the original \(D\).

Under the strict inequalities each replacement from a tree
\(D\ne T\) strictly lowers its weight. Hence \(T\) is uniquely
minimum. If \(T\) is minimum and an equality holds for a non-tree
edge \(e\), choose \(f\in P_T(e)\) attaining the path maximum.
The exchange \(T-f+e\) gives a different tree of the same
weight. These arguments also cover a tree with no non-tree edges
between distinct vertices.
\end{proof}

For the spanning tree algorithms, graphic polytope, and certificates
see \cite[Sections~50.1--50.4 and 50.6]{schrijver}.

\section{Intersection and union}

\begin{lemma}[Restrictions and contraction]
\label{p205:matroid-minors}
For a finite matroid \(M\), restriction to \(U\subseteq E\) keeps
the independent subsets of \(U\) and preserves their ranks.
For \(C\subseteq E\), contraction \(M/C\) is a matroid on
\(E\setminus C\). Its rank on \(S\subseteq E\setminus C\) is
\(r_M(S\cup C)-r_M(C)\). For any basis \(I\) of \(C\), a set
\(J\subseteq E\setminus C\) is independent in \(M/C\) exactly
when \(I\cup J\) is independent in \(M\).
\end{lemma}

\begin{proof}
Restriction inherits all independence axioms. For contraction
fix a basis \(I\) of \(C\) and take as independent those \(J\)
for which \(I\cup J\) is independent in \(M\).
This gives heredity and the empty set. If two such \(J,K\)
satisfy \(|J|<|K|\), then augment \(I\cup J\) using \(I\cup K\).
The added element lies in \(K\setminus J\). This proves the
matroid axioms on \(E\setminus C\).

Any independent \(I\cup J\) with \(J\subseteq S\) has size at
most \(r_M(C\cup S)\). For equality extend \(I\) to a maximal
independent subset of \(C\cup S\). All added elements lie in
\(S\), since no element of \(C\setminus I\) can extend the
original basis \(I\) and heredity continues to forbid it.
Thus the maximum size of \(J\) is \(r_M(C\cup S)-|I|\).
Because \(|I|=r_M(C)\), the resulting rank is independent of
the chosen basis \(I\). A set is independent exactly when its
rank equals its size, so the construction is independent of
that choice as well.
\end{proof}

This also gives the part of the minor calculus from
\cite[Chapter~10, W69]{p204} needed in the next argument.

\begin{theorem}[Edmonds' unweighted matroid intersection theorem]
\label{p205:W49}
Let \(M_1,M_2\) be finite matroids on the same ground set \(E\)
with rank functions \(r_1,r_2\). The maximum size of a common
independent set is
\[
 m=\min_{U\subseteq E}\bigl(r_1(U)+r_2(E\setminus U)\bigr).
\]
\end{theorem}

\begin{proof}
A common independent \(I\) splits into independent subsets of
\(U\) in \(M_1\) and of \(E\setminus U\) in \(M_2\). This proves
\(|I|\le m\). For attainment use induction on \(|E|\).
If \(E\) is empty both sides are zero. If \(E=\{e\}\), the
minimum is \(\min(r_1(\{e\}),r_2(\{e\}))\), which is \(1\)
exactly when \(\{e\}\) is independent in both matroids.

Suppose first that some nonempty proper \(U\subset E\) attains
the minimum, and put \(V=E\setminus U\). On \(U\) apply
induction to \(M_1|U\) and \(M_2/V\). For \(A\subseteq U\) the
corresponding rank sum is
\[
 r_1(A)+r_2(E\setminus A)-r_2(V)\ge m-r_2(V)=r_1(U),
\]
by Lemma~\ref{p205:matroid-minors}. At \(A=U\) it equals
\(r_1(U)\). Hence a common independent \(I_U\) of size \(r_1(U)\)
exists. It is a basis of \(U\) in \(M_1\) and independent in
\(M_2/V\).
On \(V\) apply induction to \(M_1/U\) and \(M_2|V\).
For \(A\subseteq V\) their rank sum is
\[
 r_1(U\cup A)-r_1(U)+r_2(V\setminus A)\ge m-r_1(U)=r_2(V),
\]
with equality at \(A=\varnothing\). This gives a common
independent \(I_V\) of size \(r_2(V)\) which is independent in
\(M_1/U\) and is a basis of \(V\) in \(M_2\).
The contraction criterion applied with basis \(I_U\) in \(M_1\)
and basis \(I_V\) in \(M_2\) makes \(I_U\cup I_V\) independent
in both. It has size \(m\).

Suppose now that every minimizing set is either \(\varnothing\)
or \(E\), where \(|E|\ge2\). Fix \(e\in E\) and put
\(E'=E\setminus\{e\}\). For any \(A\subseteq E'\) write
\(B=E'\setminus A\) and \(q=r_1(A)+r_2(B)\).
The original splits at \(A\) and \(A\cup\{e\}\) have rank
sums \(r_1(A)+r_2(B\cup\{e\})\) and
\(r_1(A\cup\{e\})+r_2(B)\), respectively.
Each is at least \(m\) and at most \(q+1\), since adding one
element increases rank by at most one. If \(q<m\), integrality
makes both original sums equal \(m\). This gives two minimizing
sets differing by one element, so they cannot both be among
\(\varnothing,E\). Thus \(q\ge m\) for every \(A\subseteq E'\).
By induction the restrictions to \(E'\) admit a common independent
set of size at least \(m\). It is common independent in the original
matroids, where the upper bound proves equality.
\end{proof}

\begin{theorem}[Matroid union, rank, and partition criterion]
\label{p205:W50}
Let \(M_1,\ldots,M_k\) be finite matroids on the same ground set
\(E\), with ranks \(r_1,\ldots,r_k\). The sets expressible as
a union of one independent set of each \(M_i\) are the independent
sets of a matroid. Its rank on \(S\subseteq E\) is
\[
 \rho(S)=\min_{T\subseteq S}
       \left(|S\setminus T|+\sum_{i=1}^k r_i(T)\right).
\]
In particular, \(S\) partitions into disjoint sets \(S_i\)
independent in the respective \(M_i\) if and only if
\[
 |T|\le \sum_{i=1}^k r_i(T)\qquad\text{for every }T\subseteq S.
\]
Empty parts are allowed. In the case \(k=0\) the only independent
set is empty and the formulas give rank zero.
\end{theorem}

\begin{proof}
If an element occurs in more than one chosen independent set,
assign it to one such index and delete it from the others.
Heredity preserves independence, so the union may be made disjoint.
For \(k=0\), \(\min_{T\subseteq S}|S\setminus T|=0\) and the
partition condition holds precisely for \(S=\varnothing\).
Assume \(k\ge1\).

Fix \(S\subseteq E\) and use the labeled ground set
\(X=S\times\{1,\ldots,k\}\). On \(X\) define two matroids.
The first allows a set \(J\subseteq X\) when each slice
\(J_i=\{e\mid(e,i)\in J\}\) is independent in \(M_i\).
The second allows at most one copy of any element \(e\).
Both are hereditary and contain the empty set.
For augmentation in the first, a larger total size yields
some slice of larger size, where the original augmentation
axiom applies. For the second, a larger set uses more distinct
labels \(e\), so an element with a new label can be added.
The rank of the first on \(U\subseteq X\) is \(\sum_i r_i(U_i)\)
by optimizing the slices separately. The rank of the second
on \(Y\subseteq X\) is the number of distinct elements \(e\)
represented in \(Y\).

A common independent set of the two matroids represents a
subset of \(S\) with its disjoint assignment to the \(M_i\).
By Theorem~\ref{p205:W49} its largest size is the minimum of
the split rank sums on \(X\). Given \(U\subseteq X\) let
\(T\subseteq S\) consist of the elements all of whose copies
belong to \(U\). The second matroid has rank \(|S\setminus T|\)
on \(X\setminus U\), and the first has rank at least
\(\sum_i r_i(T)\) on \(U\), since \(T\subseteq U_i\) for every
\(i\). Conversely \(U=T\times\{1,\ldots,k\}\) realizes equality
in this rank comparison for every \(T\subseteq S\). This proves
the formula for the largest assignable subset size \(\rho(S)\).

It remains to verify that the assignable subsets form a matroid.
Their family contains the empty set and is hereditary. The
function \(\rho\) is integral, monotone, and normalized, since
it is a maximum subset size. For \(e\notin S\), deleting \(e\)
from an assignable subset of \(S\cup\{e\}\) leaves one in \(S\),
so \(\rho(S\cup\{e\})-\rho(S)\le1\).
To show submodularity, choose minimizing \(T_1\subseteq S_1\),
\(T_2\subseteq S_2\) in the formula for \(\rho(S_1)\) and
\(\rho(S_2)\). Cardinality addition for unions and intersections
gives
\begin{align*}
 |S_1\setminus T_1|+|S_2\setminus T_2|
 &= |(S_1\cup S_2)\setminus(T_1\cup T_2)|\\
 &\quad+|(S_1\cap S_2)\setminus(T_1\cap T_2)|.
\end{align*}
Each \(r_i\) is submodular. For a direct verification, extend a
basis of \(T_1\cap T_2\) to a basis \(B\) of \(T_1\), then extend
\(B\) to a basis of \(T_1\cup T_2\). The second extension adds
only elements from \(T_2\setminus T_1\). These added elements,
together with the original basis of \(T_1\cap T_2\), form an
independent subset of \(T_2\). Its size gives
\(r_i(T_1\cap T_2)+r_i(T_1\cup T_2)-r_i(T_1)\le r_i(T_2)\).
Thus the sum of the two minima at \(T_1,T_2\) is at least the
sum of the expressions at \(T_1\cup T_2\) and \(T_1\cap T_2\)
for \(S_1\cup S_2\) and \(S_1\cap S_2\). The latter sum is
at least \(\rho(S_1\cup S_2)+\rho(S_1\cap S_2)\).

For completeness this submodularity gives augmentation of the
assignable family directly. If \(I,J\) are assignable and no
\(e\in J\setminus I\) augments \(I\), then
\(\rho(I\cup\{e\})=\rho(I)=|I|\) for each such \(e\).
For any \(A\supseteq I\) and \(e\notin A\), submodularity applied
to \(A\) and \(I\cup\{e\}\) gives
\(\rho(A\cup\{e\})-\rho(A)\le
  \rho(I\cup\{e\})-\rho(I)\).
Consequently adding the elements of \(J\setminus I\) one at a
time leaves \(\rho(I\cup J)=\rho(I)=|I|\). Monotonicity yields
\(|J|=\rho(J)\le |I|\). This proves augmentation whenever
\(|I|<|J|\), and hence \(\rho\) is the rank of the union matroid.

Finally \(S\) is assignable exactly when \(\rho(S)=|S|\).
Every expression in the minimum for \(\rho(S)\) is at least
\(|S|\) exactly when every \(T\subseteq S\) has
\(\sum_i r_i(T)\ge |T|\). Since the expression at
\(T=\varnothing\) equals \(|S|\), this is the claimed equivalence.
\end{proof}

The unweighted min--max and union formulas are in
\cite[Sections~41.1 and 42.1]{schrijver}.

\backmatter
\begin{thebibliography}{99}
\bibitem{diestel} R. Diestel, \emph{Graph Theory}, 5th ed.,
  Springer, 2017, Sections 7.1 and 9.1.
\bibitem{jukna} S. Jukna, \emph{Extremal Combinatorics}, 2nd ed.,
  Springer, 2011.
\bibitem{kst} T. K\H{o}v\'ari, V. T. S\'os, and P. Tur\'an,
  On a problem of K. Zarankiewicz, \emph{Colloquium Math.} 3
  (1954), 50--57.
  \href{https://doi.org/10.4064/cm-3-1-50-57}{DOI}.
\bibitem{zhao} Y. Zhao, \emph{Graph Theory and Additive
  Combinatorics}, MIT 18.225 lecture notes, Fall 2023,
  Sections 1.3 and 1.5.
\bibitem{alon} N. Alon and J. H. Spencer,
  \emph{The Probabilistic Method}, 4th ed., Wiley, 2016.
\bibitem{tao} T. Tao, \emph{275A, Notes 1: Integration and
  Expectation}, 2015, Exercise 44.
\bibitem{roch} S. Roch, \emph{First and Second Moment Methods},
  graduate probability lecture notes, Lecture 6.
\bibitem{vondrak} J. Vondr\'ak, \emph{The Lov\'asz Local Lemma},
  Stanford Math 233A lecture notes, 2018, Lecture 2.
\bibitem{hoeffding} W. Hoeffding, Probability inequalities for sums
  of bounded random variables, \emph{J. Amer. Statist. Assoc.} 58
  (1963), 13--30.
  \href{https://doi.org/10.1080/01621459.1963.10500830}{DOI}.
\bibitem{schrijver} A. Schrijver, \emph{Combinatorial Optimization:
  Polyhedra and Efficiency}, vols.~A--C, Springer, 2003.
\bibitem{oxley} J. Oxley, \emph{Matroid Theory}, 2nd ed.,
  Oxford University Press, 2011.
\bibitem{kuhn} H. W. Kuhn, The Hungarian method for the assignment
  problem, \emph{Naval Research Logistics Quarterly} 2 (1955), 83--97.
  \href{https://doi.org/10.1002/nav.3800020109}{DOI}.
\bibitem{p106} MathGenome, \emph{Variational Methods, Convex
  Analysis, and Optimality}, unified notes P1-06, Chapters 1 and 5,
  8 November 2026.
\bibitem{p204} MathGenome, \emph{Finite Combinatorics:
  Enumeration, Graphs, and Matroids}, unified notes P2-04,
  8 November 2026.
\bibitem{d11} MathGenome, \emph{Unified Real Analysis and Measure
  Theory Notes}, D11, Chapters 5, 6, and 10.
\bibitem{mathlib} The Mathlib Community, \emph{Mathlib},
  \href{https://github.com/leanprover-community/mathlib4/tree/56a005855ea17e0b0958f73ec801ae140b0a1679}{fixed revision
  \texttt{56a005855ea1}}. Relevant modules include
  \texttt{Combinatorics.SetFamily.KruskalKatona} and
  \texttt{Analysis.Convex.Birkhoff}.
\end{thebibliography}
\end{document}
